
\Data\ID{1003019302}
\Updated{2022-08-25 10:08:56}
\Level{A}
\SubArea{Functions with absolute values}
\LANGen{
\Question{
Which of the following functions has minimum at \( x=3 \)?}
{\( f(x)=2|x-3|+1 \)}{\( h(x)=2|x+3|-1 \)}{\( g(x)=-2|x-3|+3 \)}{\( m(x)=2|x+3| \)}{}{}}
\LANGpl{
\Question{
Która z poniższych funkcji ma minimum w  \( x=3 \)?}
{\( f(x)=2|x-3|+1 \)}{\( h(x)=2|x+3|-1 \)}{\( g(x)=-2|x-3|+3 \)}{\( m(x)=2|x+3| \)}{}{}}
\LANGcs{
\Question{
Která z následujících funkcí má minimum v bodě 3?}
{\( f(x)=2|x-3|+1 \)}{\( h(x)=2|x+3|-1 \)}{\( g(x)=-2|x-3|+3 \)}{\( m(x)=2|x+3| \)}{}{}}
\LANGsk{
\Question{
Ktorá z následujúcich funkcií má minimum v bode 3?}
{\( f(x)=2|x-3|+1 \)}{\( h(x)=2|x+3|-1 \)}{\( g(x)=-2|x-3|+3 \)}{\( m(x)=2|x+3| \)}{}{}}
\LANGes{
\Question{
¿Cuál de las siguientes funciones tiene un mínimo en \( x=3 \)?}
{\( f(x)=2|x-3|+1 \)}{\( h(x)=2|x+3|-1 \)}{\( g(x)=-2|x-3|+3 \)}{\( m(x)=2|x+3| \)}{}{}}
\End

\Data\ID{2000003901}
\Updated{2021-12-26 08:12:30}
\Level{A}
\SubArea{Functions with absolute values}
\IMG{2000003901_3-figure0.pdf}
\LANGen{
\Question{
Determine for which values of the real parameter \(a\) is the graph in the picture the graph of the function \(f(x)=|x+a|-3\); \(x \in [ -3;3]\).}
{\(a=2\)}{\(a=-2\)}{\(a=1\)}{\(a=-1\)}{}{}}
\LANGpl{
\Question{
Określ, dla których wartości rzeczywistego parametru \(a\) wykres na rysunku jest wykresem funkcji \(f(x)=|x+a|-3\); \(x \in \langle -3;3\rangle\).}
{\(a=2\)}{\(a=-2\)}{\(a=1\)}{\(a=-1\)}{}{}}
\LANGcs{
\Question{
Určete, pro které hodnoty reálného parametru \(a\) je graf na obrázku grafem funkce \(f(x)=|x+a|-3\); \(x \in \langle -3;3\rangle\).}
{\(a=2\)}{\(a=-2\)}{\(a=1\)}{\(a=-1\)}{}{}}
\LANGsk{
\Question{
Určte, pre ktoré hodnoty reálneho parametra \(a\) je graf na obrázku grafom funkcie \(f(x)=|x+a|-3\); \(x \in \langle -3;3\rangle\).}
{\(a=2\)}{\(a=-2\)}{\(a=1\)}{\(a=-1\)}{}{}}
\LANGes{
\Question{
Averigua para qué valores del parámetro real \(a\)  la gráfica en el dibujo corresponde a la función \(f(x)=|x+a|-3\); \(x \in [ -3;3]\).}
{\(a=2\)}{\(a=-2\)}{\(a=1\)}{\(a=-1\)}{}{}}
\End

\Data\ID{2000003902}
\Updated{2021-12-26 08:12:07}
\Level{A}
\SubArea{Functions with absolute values}
\IMG{2000003901_3-figure1.pdf}
\LANGen{
\Question{
Determine for which values of the real parameter \(b\) is the graph in the picture the graph of the function \(f(x)=|x-1|-b\); \(x \in [ -3;3]\).}
{\( b=-2\)}{\( b=2\)}{\( b=-1\)}{\( b=1\)}{}{}}
\LANGpl{
\Question{
Określ, dla których wartości rzeczywistego parametru \(b\) wykres na rysunku jest wykresem funkcji \(f(x)=|x-1|-b\); \(x \in \langle -3;3\rangle\).}
{\( b=-2\)}{\( b=2\)}{\( b=-1\)}{\( b=1\)}{}{}}
\LANGcs{
\Question{
Určete, pro které hodnoty reálného parametru \(b\) je graf na obrázku grafem funkce \(f(x)=|x-1|-b\); \(x \in \langle -3;3\rangle\).}
{\( b=-2\)}{\( b=2\)}{\( b=-1\)}{\( b=1\)}{}{}}
\LANGsk{
\Question{
Určte, pre ktoré hodnoty reálneho parametra \(b\) je graf na obrázku grafom funkcie \(f(x)=|x-1|-b\); \(x \in \langle -3;3\rangle\).}
{\( b=-2\)}{\( b=2\)}{\( b=-1\)}{\( b=1\)}{}{}}
\LANGes{
\Question{
Averigua para qué valores del parámetro real \(b\) la gráfica en el dibujo corresponde a la función \(f(x)=|x-1|-b\); \(x \in [ -3;3]\).}
{\( b=-2\)}{\( b=2\)}{\( b=-1\)}{\( b=1\)}{}{}}
\End

\Data\ID{2000003903}
\Updated{2021-12-26 08:12:46}
\Level{A}
\SubArea{Functions with absolute values}
\IMG{2000003901_3-figure2.pdf}
\LANGen{
\Question{
Determine for which values of the real parameters \(a\) and \(b\) is the graph in the picture the graph of the function \(f(x)=|x-a|+b\); \(x \in [ -3;4]\).}
{\( a=2\); \(b=1\)}{\( a=1\); \(b=2\)}{\( a=-2\); \(b=1\)}{\( a=2\); \(b=-1\)}{}{}}
\LANGpl{
\Question{
Określ, dla których wartości rzeczywistych parametrów \(a\) i \(b\) wykres na rysunku jest wykresem funkcji \(f(x)=|x-a|+b\); \(x \in \langle -3;4\rangle\).}
{\( a=2\); \(b=1\)}{\( a=1\); \(b=2\)}{\( a=-2\); \(b=1\)}{\( a=2\); \(b=-1\)}{}{}}
\LANGcs{
\Question{
Určete, pro které hodnoty reálných parametrů \(a\) a \(b\) je graf na obrázku grafem funkce \(f(x)=|x-a|+b\); \(x \in \langle -3;4\rangle\).}
{\( a=2\); \(b=1\)}{\( a=1\); \(b=2\)}{\( a=-2\); \(b=1\)}{\( a=2\); \(b=-1\)}{}{}}
\LANGsk{
\Question{
Určte, pre ktoré hodnoty reálnych parametrov \(a\) a \(b\) je graf na obrázku grafom funkcie \(f(x)=|x-a|+b\); \(x \in \langle -3;4\rangle\).}
{\( a=2\); \(b=1\)}{\( a=1\); \(b=2\)}{\( a=-2\); \(b=1\)}{\( a=2\); \(b=-1\)}{}{}}
\LANGes{
\Question{
Averigua para qué valores de los parámetros reales \(a\) y \(b\) la gráfica del dibujo corresponde a la función \(f(x)=|x-a|+b\); \(x \in [ -3;4]\).}
{\( a=2\); \(b=1\)}{\( a=1\); \(b=2\)}{\( a=-2\); \(b=1\)}{\( a=2\); \(b=-1\)}{}{}}
\End

\Data\ID{2110009101}
\Updated{2022-11-02 09:11:56}
\Level{A}
\SubArea{Functions with absolute values}
\LANGen{
\Question{
Identify a possible graph of the function
\(f(x)=2 -|x| \).}
{}{}{}{}{}{}}
\LANGpl{
\Question{
Wskaż wykres funkcji
\(f(x)=2 -|x| \).}
{}{}{}{}{}{}}
\LANGcs{
\Question{
Který z uvedených grafů je grafem funkce
\(f(x)=2 -|x| \)?}
{}{}{}{}{}{}}
\LANGsk{
\Question{
Ktorý z grafov je grafom funkcie \(f(x)=2 -|x| \)?}
{}{}{}{}{}{}}
\LANGes{
\Question{
Identifica la gráfica de la función
\(f(x)=2 -|x| \).}
{}{}{}{}{}{}}

\IMGanswer{2010009101-figure0.pdf}
\IMGanswer{2010009101-figure1.pdf}
\IMGanswer{2010009101-figure2.pdf}
\IMGanswer{2010009101-figure3.pdf}\End

\Data\ID{2110009102}
\Updated{2022-11-02 09:11:56}
\Level{A}
\SubArea{Functions with absolute values}
\LANGen{
\Question{
Identify a possible graph of the function
\(f(x)= -|x-2| \).}
{}{}{}{}{}{}}
\LANGpl{
\Question{
Wskaż wykres funkcji 
\(f(x)= -|x-2| \).}
{}{}{}{}{}{}}
\LANGcs{
\Question{
Který z grafů je grafem funkce
\(f(x)= -|x-2| \)?}
{}{}{}{}{}{}}
\LANGsk{
\Question{
Ktorý z grafov je grafom funkcie \(f(x)= -|x-2| \)?}
{}{}{}{}{}{}}
\LANGes{
\Question{
Identifica la gráfica de la función
\(f(x)= -|x-2| \).}
{}{}{}{}{}{}}

\IMGanswer{2010009101-figure4.pdf}
\IMGanswer{2010009101-figure5.pdf}
\IMGanswer{2010009101-figure6.pdf}
\IMGanswer{2010009101-figure7.pdf}\End

\Data\ID{2110009103}
\Updated{2022-11-02 09:11:56}
\Level{A}
\SubArea{Functions with absolute values}
\LANGen{
\Question{
Identify a possible graph of the function
\(f(x)= |x| +\sqrt5\).}
{}{}{}{}{}{}}
\LANGpl{
\Question{
Wskaż wykres funkcji
\(f(x)= |x| +\sqrt5\).}
{}{}{}{}{}{}}
\LANGcs{
\Question{
Který z grafů je grafem funkce
\(f(x)= |x| +\sqrt5\)?}
{}{}{}{}{}{}}
\LANGsk{
\Question{
Ktorý z nasledujúcich grafov je grafom funkcie \(f(x)= |x| +\sqrt5\)?}
{}{}{}{}{}{}}
\LANGes{
\Question{
Identifica la gráfica de la función
\(f(x)= |x| +\sqrt5\).}
{}{}{}{}{}{}}

\IMGanswer{2010009101-figure8.pdf}
\IMGanswer{2010009101-figure9.pdf}
\IMGanswer{2010009101-figure10.pdf}
\IMGanswer{2010009101-figure11.pdf}\End

\Data\ID{2110009208}
\Updated{2022-11-02 09:11:56}
\Level{A}
\SubArea{Functions with absolute values}
\LANGen{
\Question{
Identify a possible graph of the function
\(f(x)= |2+x|\).}
{}{}{}{}{}{}}
\LANGpl{
\Question{
Wskaż wykres funkcji
\(f(x)= |2+x|\).}
{}{}{}{}{}{}}
\LANGcs{
\Question{
Který z grafů je grafem funkce
\(f(x)= |2+x|\)?}
{}{}{}{}{}{}}
\LANGsk{
\Question{
Ktorý z nasledujúcich grafov je grafom funkcie \(f(x)= |2+x|\)?}
{}{}{}{}{}{}}
\LANGes{
\Question{
Identifica la gráfica de la función
\(f(x)= |2+x|\).}
{}{}{}{}{}{}}

\IMGanswer{2010009201-figure8_2.pdf}
\IMGanswer{2010009201-figure9_2.pdf}
\IMGanswer{2010009201-figure10_2.pdf}
\IMGanswer{2010009201-figure11_2.pdf}\End

\Data\ID{9000024403}
\Updated{2022-04-23 05:04:37}
\Level{A}
\SubArea{Functions with absolute values}
\IMG{000244_03-figure0.pdf}
\LANGen{
\Question{
Identify the value of a real parameter \(a\) 
such that the graph of the function 
\[ 
 f(x) = |x - a|- 3
\]
corresponds to the picture.}
{\(\ \ 1\)}{\(\ \ - 2\)}{\(\ \ - 3\)}{\(\ \ 4\)}{\(\ \ 3\)}{}}
\LANGpl{
\Question{
Określ wartość rzeczywistego parametru \(a\), dla którego wykres funkcji 
\[ 
 f\colon y = |x - a|- 3
\]
przedstawiony jest na rysunku poniżej.}
{\(\ \ 1\)}{\(\ \ - 2\)}{\(\ \ - 3\)}{\(\ \ 4\)}{\(\ \ 3\)}{}}
\LANGcs{
\Question{
Určete hodnotu reálného čísla \(a\) 
tak, aby na obrázku byl graf funkce dané předpisem
\(f\colon y = |x - a|- 3\).}
{\(\ \ 1\)}{\(\ \ - 2\)}{\(\ \ - 3\)}{\(\ \ 4\)}{\(\ \ 3\)}{}}
\LANGsk{
\Question{
Rozhodnite, pre ktoré hodnoty reálneho parametra \(a\) 
je daný graf na obrázku grafom funkcie určenej predpisom
\(f\colon y = |x - a|- 3\).}
{\(\ \ 1\)}{\(\ \ - 2\)}{\(\ \ - 3\)}{\(\ \ 4\)}{\(\ \ 3\)}{}}
\LANGes{
\Question{
Identifica el valor del parámetro real \(a\) 
tal que la gráfica de la función
\[ 
 f(x) = |x - a|- 3
\]
corresponda al dibujo.}
{\(\ \ 1\)}{\(\ \ - 2\)}{\(\ \ - 3\)}{\(\ \ 4\)}{\(\ \ 3\)}{}}
\End

\Data\ID{9000024404}
\Updated{2020-11-24 11:11:10}
\Level{A}
\SubArea{Functions with absolute values}
\IMG{000244_04-figure0.pdf}
\LANGen{
\Question{
Identify the value of a real parameter \(b\) 
such that the graph of the function 
\[ 
 f(x) = |x - 0.25| + b
\]
corresponds to the picture.}
{\(\ \ 0.5\)}{\(\ \ - 0.5\)}{\(\ \ 0.75\)}{\(\ \ - 0.75\)}{\(\ \ 0.25\)}{}}
\LANGpl{
\Question{
Określ wartość rzeczywistego parametru \(b\), dla którego wykres funkcji 
\[ 
 f\colon y = |x - 0.25| + b
\]
przedstawiony jest na rysunku poniżej.}
{\(\ \ 0.5\)}{\(\ \ - 0.5\)}{\(\ \ 0.75\)}{\(\ \ - 0.75\)}{\(\ \ 0.25\)}{}}
\LANGcs{
\Question{
Uvažujte funkci \(f\) 
zadanou grafem na obrázku. Zjistěte, pro jakou hodnotu reálného čísla
\(b\) platí,
že \(f\colon y = |x - 0{,}25| + b\).}
{\(\ \ 0{,}5\)}{\(\ \ - 0{,}5\)}{\(\ \ 0{,}75\)}{\(\ \ - 0{,}75\)}{\(\ \ 0{,}25\)}{}}
\LANGsk{
\Question{
Rozhodnite, pre ktoré hodnoty reálneho parametra \(b\) je daný graf na obrázku grafom funkcie určenej predpisom \(f\colon y = |x - 0{,}25| + b\).}
{\(\ \ 0{,}5\)}{\(\ \ - 0{,}5\)}{\(\ \ 0{,}75\)}{\(\ \ - 0{,}75\)}{\(\ \ 0{,}25\)}{}}
\LANGes{
\Question{
Identifica el valor del parámetro real \(b\) 
tal que la gráfica de la función 
\[ 
 f(x) = |x - 0.25| + b
\]
corresponda al dibujo.}
{\(\ \ 0.5\)}{\(\ \ - 0.5\)}{\(\ \ 0.75\)}{\(\ \ - 0.75\)}{\(\ \ 0.25\)}{}}
\End

\Data\ID{9000024405}
\Updated{2020-11-24 11:11:23}
\Level{A}
\SubArea{Functions with absolute values}
\IMG{000244_05-figure0.pdf}
\LANGen{
\Question{
Identify the value of a real parameter \(b\) 
such that the graph of the function 
\[ 
 f(x) = \left |x -\frac{1} 
{3}\right | + b
\]
corresponds to the picture.}
{\(\ \ - 1\)}{\(\ \ -\frac{2} 
{3}\)}{\(\ \ \frac{2}
{3}\)}{\(\ \ \frac{1}
{3}\)}{\(\ \ 1\)}{}}
\LANGpl{
\Question{
Określ wartość rzeczywistego parametru \(b\), dla którego wykres funkcji
\[ 
 f\colon y = \left |x -\frac{1} 
{3}\right | + b
\]
przedstawiony jest na rysunku poniżej.}
{\(\ \ - 1\)}{\(\ \ -\frac{2} 
{3}\)}{\(\ \ \frac{2}
{3}\)}{\(\ \ \frac{1}
{3}\)}{\(\ \ 1\)}{}}
\LANGcs{
\Question{
Uvažujte funkci \(f\) 
zadanou grafem na obrázku. Zjistěte, pro jakou hodnotu reálného čísla
\(b\) platí,
že \(f\colon y = \left |x -\frac{1} 
{3}\right | + b\).}
{\(\ \ - 1\)}{\(\ \ -\frac{2} 
{3}\)}{\(\ \ \frac{2}
{3}\)}{\(\ \ \frac{1}
{3}\)}{\(\ \ 1\)}{}}
\LANGsk{
\Question{
Rozhodnite, pre ktoré hodnoty reálneho parametra \(b\) je daný graf na obrázku grafom funkcie určenej predpisom  \(f\colon y = \left |x -\frac{1} 
{3}\right | + b\).}
{\(\ \ - 1\)}{\(\ \ -\frac{2} 
{3}\)}{\(\ \ \frac{2}
{3}\)}{\(\ \ \frac{1}
{3}\)}{\(\ \ 1\)}{}}
\LANGes{
\Question{
Identifica el valor del parámetro real \(b\) 
tal que la gráfica de la función 
\[ 
 f(x) = \left |x -\frac{1} 
{3}\right | + b
\]
corresponda al dibujo.}
{\(\ \ - 1\)}{\(\ \ -\frac{2} 
{3}\)}{\(\ \ \frac{2}
{3}\)}{\(\ \ \frac{1}
{3}\)}{\(\ \ 1\)}{}}
\End

\Data\ID{9000024406}
\Updated{2022-04-23 05:04:37}
\Level{A}
\SubArea{Functions with absolute values}
\IMG{000244_06-figure0.pdf}
\LANGen{
\Question{
Identify the values of real parameters \(a\) 
and \(b\) 
such that the graph of the function 
\[ 
 f(x)= |x + a| + b
\]
corresponds to the picture.}
{\(\ \ a = 3,\quad \phantom{ -} b = 2\)}{\(\ \ a = 2,\quad \phantom{ -} b = 3\)}{\(\ \ a = 2,\quad \phantom{ -} b = -3\)}{\(\ \ a = -3,\quad b = 2\)}{}{}}
\LANGpl{
\Question{
Określ wartości rzeczywiste parametrów \(a\) 
i \(b\), dla których wykres funkcji
\[ 
 f\colon y = |x + a| + b
\]
przedstawiony jest na rysunku poniżej.}
{\(\ \ a = 3,\quad \phantom{ -} b = 2\)}{\(\ \ a = 2,\quad \phantom{ -} b = 3\)}{\(\ \ a = 2,\quad \phantom{ -} b = -3\)}{\(\ \ a = -3,\quad b = 2\)}{}{}}
\LANGcs{
\Question{
Uvažujte funkci \(f\) 
zadanou grafem na obrázku. Zjistěte, pro jaké hodnoty reálných čísel
\(a\),
\(b\) platí,
že \(f\colon y = |x + a| + b\).}
{\(\ \ a = 3,\quad \phantom{ -} b = 2\)}{\(\ \ a = 2,\quad \phantom{ -} b = 3\)}{\(\ \ a = 2,\quad \phantom{ -} b = -3\)}{\(\ \ a = -3,\quad b = 2\)}{}{}}
\LANGsk{
\Question{
Rozhodnite, pre ktoré hodnoty reálnych parametrov \(a\),
\(b\) je daný graf na obrázku grafom funkcie určenej predpisom \(f\colon y = |x + a| + b\).}
{\(\ \ a = 3,\quad \phantom{ -} b = 2\)}{\(\ \ a = 2,\quad \phantom{ -} b = 3\)}{\(\ \ a = 2,\quad \phantom{ -} b = -3\)}{\(\ \ a = -3,\quad b = 2\)}{}{}}
\LANGes{
\Question{
Identifica los valores de los parámetros reales \(a\) 
y \(b\) 
tales que la gráfica de la función 
\[ 
 f(x)= |x + a| + b
\]
corresponda al dibujo.}
{\(\ \ a = 3,\quad \phantom{ -} b = 2\)}{\(\ \ a = 2,\quad \phantom{ -} b = 3\)}{\(\ \ a = 2,\quad \phantom{ -} b = -3\)}{\(\ \ a = -3,\quad b = 2\)}{}{}}
\End

\Data\ID{9000024407}
\Updated{2021-04-11 09:04:58}
\Level{A}
\SubArea{Functions with absolute values}
\IMG{000244_07-figure0.pdf}
\LANGen{
\Question{
Identify the value of a real parameter \(a\) 
such that the graph of the function 
\[ 
 f(x) = |x + a|- 2
\]
corresponds to the picture.}
{\(\ \ - 1\)}{\(\ \ 1\)}{\(\ \ - 2\)}{\(\ \ 2\)}{\(\ \ 3\)}{}}
\LANGpl{
\Question{
Określ wartość rzeczywistego parametru \(a\), dla którego wykres funkcji
\[ 
 f\colon y = |x + a|- 2
\]
przedstawiony jest na rysunku poniżej.}
{\(\ \ - 1\)}{\(\ \ 1\)}{\(\ \ - 2\)}{\(\ \ 2\)}{\(\ \ 3\)}{}}
\LANGcs{
\Question{
Určete hodnotu reálného čísla
\(a\) tak, aby na obrázku
byl graf funkce \(f\colon y = |x + a|- 2\).}
{\(\ \ - 1\)}{\(\ \ 1\)}{\(\ \ - 2\)}{\(\ \ 2\)}{\(\ \ 3\)}{}}
\LANGsk{
\Question{
Rozhodnite, pre ktoré hodnoty reálneho parametra \(a\) je daný graf na obrázku grafom funkcie určenej predpisom \(f\colon y = |x + a|- 2\).}
{\(\ \ - 1\)}{\(\ \ 1\)}{\(\ \ - 2\)}{\(\ \ 2\)}{\(\ \ 3\)}{}}
\LANGes{
\Question{
Identifica el valor del parámetro real \(a\) 
tal que la gráfica de la función 
\[ 
 f(x) = |x + a|- 2
\]
corresponda al dibujo.}
{\(\ \ - 1\)}{\(\ \ 1\)}{\(\ \ - 2\)}{\(\ \ 2\)}{\(\ \ 3\)}{}}
\End

\Data\ID{9000024408}
\Updated{2020-11-24 11:11:27}
\Level{A}
\SubArea{Functions with absolute values}
\IMG{000244_08-figure0.pdf}
\LANGen{
\Question{
Identify the values of real parameters \(a\) 
and \(b\) 
such that the graph of the function 
\[ 
 f(x)= |x - a| + b
\]
corresponds to the picture.}
{\(\ \ a = 3,\quad \phantom{ -} b = -2\)}{\(\ \ a = -3,\quad b = 2\)}{\(\ \ a = 2,\quad \phantom{ -} b = -3\)}{\(\ \ a = -2,\quad b = 2\)}{}{}}
\LANGpl{
\Question{
Określ wartość rzeczywistych parametrów \(a\) 
i \(b\), dla których wykres funkcji
\[ 
 f\colon y = |x - a| + b
\]
przedstawiony jest na rysunku poniżej.}
{\(\ \ a = 3,\quad \phantom{ -} b = -2\)}{\(\ \ a = -3,\quad b = 2\)}{\(\ \ a = 2,\quad \phantom{ -} b = -3\)}{\(\ \ a = -2,\quad b = 2\)}{}{}}
\LANGcs{
\Question{
Uvažujte funkci \(f\) 
zadanou grafem na obrázku. Zjistěte, pro jaké hodnoty reálných čísel
\(a\),
\(b\) platí,
že \(f\colon y = |x - a| + b\).}
{\(\ \ a = 3,\quad \phantom{ -} b = -2\)}{\(\ \ a = -3,\quad b = 2\)}{\(\ \ a = 2,\quad \phantom{ -} b = -3\)}{\(\ \ a = -2,\quad b = 2\)}{}{}}
\LANGsk{
\Question{
Rozhodnite, pre ktoré hodnoty reálnych parametrov \(a\), \(b\) je daný graf na obrázku grafom funkcie určenej predpisom \(f\colon y = |x - a| + b\).}
{\(\ \ a = 3,\quad \phantom{ -} b = -2\)}{\(\ \ a = -3,\quad b = 2\)}{\(\ \ a = 2,\quad \phantom{ -} b = -3\)}{\(\ \ a = -2,\quad b = 2\)}{}{}}
\LANGes{
\Question{
Identifica los valores de los parámetros reales \(a\) 
y \(b\) 
tales que la gráfica de la función
\[ 
 f(x)= |x - a| + b
\]
corresponda al dibujo.}
{\(\ \ a = 3,\quad \phantom{ -} b = -2\)}{\(\ \ a = -3,\quad b = 2\)}{\(\ \ a = 2,\quad \phantom{ -} b = -3\)}{\(\ \ a = -2,\quad b = 2\)}{}{}}
\End

\Data\ID{9000024503}
\Updated{2020-11-24 11:11:00}
\Level{A}
\SubArea{Functions with absolute values}
\IMG{000245_03-figure0.pdf}
\LANGen{
\Question{
Identify the value of a real parameter \(a\) 
such that the graph of the function 
\[ 
 f(x) = |x - a| + 3
\]
corresponds to the picture.}
{\(1\)}{\(3\)}{\(- 3\)}{\(- 1\)}{\(4\)}{}}
\LANGpl{
\Question{
Określ wartość rzeczywistego parametru \(a\), dla którego wykres funkcji
\[ 
 f\colon y = |x - a| + 3
\]
przedstawiony jest na rysunku poniżej.}
{\(1\)}{\(3\)}{\(- 3\)}{\(- 1\)}{\(4\)}{}}
\LANGcs{
\Question{
Na obrázku je graf funkce \(f\). Zjistěte, pro jakou hodnotu reálného čísla
\(a\) platí,
že \(f\colon y = |x - a| + 3\).}
{\(1\)}{\(3\)}{\(- 3\)}{\(- 1\)}{\(4\)}{}}
\LANGsk{
\Question{
Rozhodnite, pre ktoré hodnoty reálneho parametra \(a\) je daný graf na obrázku grafom funkcie určenej predpisom \(f\colon y = |x - a| + 3\).}
{\(1\)}{\(3\)}{\(- 3\)}{\(- 1\)}{\(4\)}{}}
\LANGes{
\Question{
Identifica el valor del parámetro real \(a\) 
tal que la gráfica de la función
\[ 
 f(x) = |x - a| + 3
\]
corresponda al dibujo.}
{\(1\)}{\(3\)}{\(- 3\)}{\(- 1\)}{\(4\)}{}}
\End

\Data\ID{9000024504}
\Updated{2020-11-24 11:11:14}
\Level{A}
\SubArea{Functions with absolute values}
\IMG{000245_04-figure0.pdf}
\LANGen{
\Question{
Identify the value of a real parameter \(b\) 
such that the graph of the function 
\[ 
 f(x) = |x - 0.5| + b
\]
corresponds to the picture.}
{\(1\)}{\(- 1\)}{\(- 0.5\)}{\(0.5\)}{\(1.5\)}{}}
\LANGpl{
\Question{
Określ wartość rzeczywistego parametru  \(b\), dla którego wykres funkcji
\[ 
 f\colon y = |x - 0.5| + b
\]
przedstawiony jest na rysunku poniżej.}
{\(1\)}{\(- 1\)}{\(- 0.5\)}{\(0.5\)}{\(1.5\)}{}}
\LANGcs{
\Question{
Uvažujte funkci \(f\) 
zadanou grafem. Zjistěte, pro jakou hodnotu reálného čísla
\(b\) platí,
že \(f\colon y = |x - 0{,}5| + b\).}
{\(1\)}{\(- 1\)}{\(- 0{,}5\)}{\(0{,}5\)}{\(1{,}5\)}{}}
\LANGsk{
\Question{
Rozhodnite, pre ktoré hodnoty reálneho parametra \(b\) je daný graf na obrázku grafom funkcie určenej predpisom \(f\colon y = |x - 0{,}5| + b\).}
{\(1\)}{\(- 1\)}{\(- 0{,}5\)}{\(0{,}5\)}{\(1{,}5\)}{}}
\LANGes{
\Question{
Identifica el valor del parámetro real \(b\) 
tal que la gráfica de la función
\[ 
 f(x) = |x - 0.5| + b
\]
corresponda al dibujo.}
{\(1\)}{\(- 1\)}{\(- 0.5\)}{\(0.5\)}{\(1.5\)}{}}
\End

\Data\ID{9000024505}
\Updated{2020-11-24 11:11:42}
\Level{A}
\SubArea{Functions with absolute values}
\IMG{000245_05-figure0.pdf}
\LANGen{
\Question{
Identify the value of a real parameter \(b\) 
such that the graph of the function 
\[ 
 f(x) = \left |x + \frac{1} 
{2}\right | + b
\]
corresponds to the picture.}
{\(- 1\)}{\(-\frac{3} 
{2}\)}{\(-\frac{1} 
{2}\)}{\(\frac{1}
{2}\)}{\(1\)}{}}
\LANGpl{
\Question{
Określ wartość rzeczywistego parametru \(b\), dla którego wykres funkcji 
\[ 
 f\colon y = \left |x + \frac{1} 
{2}\right | + b
\]
przedstawiony jest na rysunku poniżej.}
{\(- 1\)}{\(-\frac{3} 
{2}\)}{\(-\frac{1} 
{2}\)}{\(\frac{1}
{2}\)}{\(1\)}{}}
\LANGcs{
\Question{
Uvažujte funkci \(f\) 
zadanou grafem. Zjistěte, pro jakou hodnotu reálného čísla
\(b\) platí,
že \(f\colon y = \left |x + \frac{1} 
{2}\right | + b\).}
{\(- 1\)}{\(-\frac{3} 
{2}\)}{\(-\frac{1} 
{2}\)}{\(\frac{1}
{2}\)}{\(1\)}{}}
\LANGsk{
\Question{
Rozhodnite, pre ktoré hodnoty reálneho parametra \(b\) je daný graf na obrázku grafom funkcie určenej predpisom \(f\colon y = \left |x + \frac{1} 
{2}\right | + b\).}
{\(- 1\)}{\(-\frac{3} 
{2}\)}{\(-\frac{1} 
{2}\)}{\(\frac{1}
{2}\)}{\(1\)}{}}
\LANGes{
\Question{
Identifica el valor del parámetro real \(b\) 
tal que la gráfica de la función
\[ 
 f(x) = \left |x + \frac{1} 
{2}\right | + b
\]
corresponda al dibujo.}
{\(- 1\)}{\(-\frac{3} 
{2}\)}{\(-\frac{1} 
{2}\)}{\(\frac{1}
{2}\)}{\(1\)}{}}
\End

\Data\ID{9000024506}
\Updated{2020-11-24 11:11:49}
\Level{A}
\SubArea{Functions with absolute values}
\IMG{000245_06-figure0.pdf}
\LANGen{
\Question{
Identify the value of real parameters \(a\) 
and \(b\) 
such that the graph of the function 
\[ 
 f(x)= \left |x + a\right | + b
\]
corresponds to the picture.}
{\(a = 2,\ b = -3\)}{\(a = 2,\ b = 3\)}{\(a = 3,\ b = 2\)}{\(a = -3,\ b = -2\)}{}{}}
\LANGpl{
\Question{
Określ wartość rzeczywistych parametrów  \(a\) 
i \(b\), dla których wykres funkcji 
\[ 
 f\colon y = \left |x + a\right | + b
\]
przedstawiony jest na rysunku poniżej.}
{\(a = 2,\ b = -3\)}{\(a = 2,\ b = 3\)}{\(a = 3,\ b = 2\)}{\(a = -3,\ b = -2\)}{}{}}
\LANGcs{
\Question{
Uvažujte funkci \(f\) 
zadanou grafem. Zjistěte, pro jaké hodnoty reálných čísel
\(a\),
\(b\) platí,
že \(f\colon y = \left |x + a\right | + b\).}
{\(a = 2,\ b = -3\)}{\(a = 2,\ b = 3\)}{\(a = 3,\ b = 2\)}{\(a = -3,\ b = -2\)}{}{}}
\LANGsk{
\Question{
Rozhodnite, pre ktoré hodnoty reálnych parametrov \(a\), \(b\) je daný graf na obrázku grafom funkcie určenej predpisom \(f\colon y = \left |x + a\right | + b\).}
{\(a = 2,\ b = -3\)}{\(a = 2,\ b = 3\)}{\(a = 3,\ b = 2\)}{\(a = -3,\ b = -2\)}{}{}}
\LANGes{
\Question{
Identifica los valores de los parámetros reales \(a\) y \(b\) 
tales que la gráfica de la función
\[ 
 f(x)= \left |x + a\right | + b
\]
corresponda al dibujo.}
{\(a = 2,\ b = -3\)}{\(a = 2,\ b = 3\)}{\(a = 3,\ b = 2\)}{\(a = -3,\ b = -2\)}{}{}}
\End

\Data\ID{9000024507}
\Updated{2020-11-24 11:11:34}
\Level{A}
\SubArea{Functions with absolute values}
\IMG{000245_07-figure0.pdf}
\LANGen{
\Question{
Identify the value of a real parameter \(a\) 
such that the graph of the function 
\[ 
 f(x) = -|a - x| + 2
\]
corresponds to the picture.}
{\(3\)}{\(2\)}{\(1\)}{\(- 1\)}{\(4\)}{}}
\LANGpl{
\Question{
Określ wartość rzeczywistego parametru \(a\), dla którego wykres funkcji 
\[ 
 f\colon y = -|a - x| + 2
\]
przedstawiony jest na rysunku poniżej.}
{\(3\)}{\(2\)}{\(1\)}{\(- 1\)}{\(4\)}{}}
\LANGcs{
\Question{
Na obrázku je graf funkce \(f\). Určete, pro jaké reálné číslo \(a\) 
je funkce \(f\) dána
předpisem \(y = -|a - x| + 2\).}
{\(3\)}{\(2\)}{\(1\)}{\(- 1\)}{\(4\)}{}}
\LANGsk{
\Question{
Rozhodnite, pre ktoré hodnoty reálneho parametra \(a\) je daný graf na obrázku grafom funkcie určenej predpisom \(f\colon y = -|a - x| + 2\).}
{\(3\)}{\(2\)}{\(1\)}{\(- 1\)}{\(4\)}{}}
\LANGes{
\Question{
Identifica el valor del parámetro real \(a\) 
tal que la gráfica de la función
\[ 
 f(x) = -|a - x| + 2
\]
corresponda al dibujo.}
{\(3\)}{\(2\)}{\(1\)}{\(- 1\)}{\(4\)}{}}
\End

\Data\ID{9000024508}
\Updated{2020-11-24 02:11:44}
\Level{A}
\SubArea{Functions with absolute values}
\IMG{000245_08-figure0.pdf}
\LANGen{
\Question{
Identify the value of real parameters \(a\) 
and \(b\) 
such that the graph of the function 
\[ 
 f(x) = \left |x - a\right | + b
\]
corresponds to the picture.}
{\(a = 2,\ b = -2\)}{\(a = -2,\ b = 2\)}{\(a = 2,\ b = 2\)}{\(a = -2,\ b = -2\)}{}{}}
\LANGpl{
\Question{
Określ wartość rzeczywistych parametrów \(a\) 
i \(b\), dla których wykres funkcji
\[ 
 f\colon y = \left |x - a\right | + b
\]
przedstawiony jest na rysunku poniżej.}
{\(a = 2,\ b = -2\)}{\(a = -2,\ b = 2\)}{\(a = 2,\ b = 2\)}{\(a = -2,\ b = -2\)}{}{}}
\LANGcs{
\Question{
Uvažujte funkci \(f\) 
zadanou grafem. Zjistěte, pro jaké hodnoty reálných čísel
\(a\),
\(b\) platí,
že \(f\colon y = \left |x - a\right | + b\).}
{\(a = 2,\ b = -2\)}{\(a = -2,\ b = 2\)}{\(a = 2,\ b = 2\)}{\(a = -2,\ b = -2\)}{}{}}
\LANGsk{
\Question{
Rozhodnite, pre ktoré hodnoty reálnych parametrov \(a\), \(b\) je daný graf na obrázku grafom funkcie určenej predpisom \(f\colon y = \left |x - a\right | + b\).}
{\(a = 2,\ b = -2\)}{\(a = -2,\ b = 2\)}{\(a = 2,\ b = 2\)}{\(a = -2,\ b = -2\)}{}{}}
\LANGes{
\Question{
Identifica el valor de los parámetros reales \(a\) 
y \(b\) 
tales que la gráfica de la función
\[ 
 f(x) = \left |x - a\right | + b
\]
corresponda al dibujo.}
{\(a = 2,\ b = -2\)}{\(a = -2,\ b = 2\)}{\(a = 2,\ b = 2\)}{\(a = -2,\ b = -2\)}{}{}}
\End

\Data\ID{9100004001}
\Updated{2022-05-29 12:05:45}
\Level{A}
\SubArea{Functions with absolute values}
\LANGen{
\Question{
Identify a possible graph of the function
\(g(x)= |x|- 1\).}
{}{}{}{}{}{}}
\LANGpl{
\Question{
Określ, który z poniższych rysunków przedstawia wykres funkcji
\(g\colon y = |x|- 1\).}
{}{}{}{}{}{}}
\LANGcs{
\Question{
Určete, na kterém z obrázků je graf funkce
\(g\colon y = |x|- 1\).}
{}{}{}{}{}{}}
\LANGsk{
\Question{
Ktorý z nasledujúcich grafov je grafom funkcie
\(g\colon y = |x|- 1\)?}
{}{}{}{}{}{}}
\LANGes{
\Question{
Identifica la posible gráfica de la función
\(g(x)= |x|- 1\).}
{}{}{}{}{}{}}

\IMGanswer{000040_01-figure2.pdf}
\IMGanswer{000040_01-figure0.pdf}
\IMGanswer{000040_01-figure1.pdf}
\IMGanswer{000040_01-figure3.pdf}\End

\Data\ID{9100004002}
\Updated{2022-05-29 12:05:45}
\Level{A}
\SubArea{Functions with absolute values}
\LANGen{
\Question{
Identify a possible graph of the function
\(g(x) = |x - 2|\).}
{}{}{}{}{}{}}
\LANGpl{
\Question{
Określ, który z poniższych rysunków przedstawia wykres funkcji
\(g\colon y = |x - 2|\).}
{}{}{}{}{}{}}
\LANGcs{
\Question{
Určete, na kterém z obrázků je graf funkce
\(g\colon y = |x - 2|\).}
{}{}{}{}{}{}}
\LANGsk{
\Question{
Ktorý z nasledujúcich grafov je grafom funkcie
\(g\colon y = |x - 2|\)?}
{}{}{}{}{}{}}
\LANGes{
\Question{
Identifica la posible gráfica de la función
\(g(x) = |x - 2|\).}
{}{}{}{}{}{}}

\IMGanswer{000040_02-figure2.pdf}
\IMGanswer{000040_02-figure0.pdf}
\IMGanswer{000040_02-figure1.pdf}
\IMGanswer{000040_02-figure3.pdf}\End

\Data\ID{9100004003}
\Updated{2022-05-29 12:05:45}
\Level{A}
\SubArea{Functions with absolute values}
\LANGen{
\Question{
Identify a possible graph of the function
\(g(x)= -|x| + 1\).}
{}{}{}{}{}{}}
\LANGpl{
\Question{
Określ, który z rysunków przedstawia wykres funkcji
\(g\colon y = -|x| + 1\).}
{}{}{}{}{}{}}
\LANGcs{
\Question{
Určete, který z obrázků je grafem funkce
\(g\colon y = -|x| + 1\).}
{}{}{}{}{}{}}
\LANGsk{
\Question{
Ktorý z grafov je grafom funkcie
\(g\colon y = -|x| + 1\)?}
{}{}{}{}{}{}}
\LANGes{
\Question{
Identifica la posible gráfica de la función
\(g(x)= -|x| + 1\).}
{}{}{}{}{}{}}

\IMGanswer{000040_03-figure3.pdf}
\IMGanswer{000040_03-figure0.pdf}
\IMGanswer{000040_03-figure1.pdf}
\IMGanswer{000040_03-figure2.pdf}\End

\Data\ID{9100004004}
\Updated{2022-05-29 12:05:45}
\Level{A}
\SubArea{Functions with absolute values}
\LANGen{
\Question{
Identify a possible graph of the function
\(g(x)= -|x + 1|\).}
{}{}{}{}{}{}}
\LANGpl{
\Question{
Określ, który z rysunków przedstawia wykres funkcji
\(g\colon y = -|x + 1|\).}
{}{}{}{}{}{}}
\LANGcs{
\Question{
Určete, který z obrázků je grafem funkce
\(g\colon y = -|x + 1|\).}
{}{}{}{}{}{}}
\LANGsk{
\Question{
Ktorý z grafov je grafom funkcie
\(g\colon y = -|x + 1|\)?}
{}{}{}{}{}{}}
\LANGes{
\Question{
Identifica la posible gráfica de la función
\(g(x)= -|x + 1|\).}
{}{}{}{}{}{}}

\IMGanswer{000040_04-figure1.pdf}
\IMGanswer{000040_04-figure0.pdf}
\IMGanswer{000040_04-figure2.pdf}
\IMGanswer{000040_04-figure3.pdf}\End

\Data\ID{9100004005}
\Updated{2020-11-24 02:11:23}
\Level{A}
\SubArea{Functions with absolute values}
\LANGen{
\Question{
Identify a possible graph of the function
\(g(x) = |x + 1|- 2\).}
{}{}{}{}{}{}}
\LANGpl{
\Question{
Określ, który z poniższych rysunków przedstawia wykres funkcji
\(g\colon y = |x + 1|- 2\).}
{}{}{}{}{}{}}
\LANGcs{
\Question{
Určete, na kterém z obrázků je graf funkce
\(g\colon y = |x + 1|- 2\).}
{}{}{}{}{}{}}
\LANGsk{
\Question{
Ktorý z nasledujúcich grafov je grafom funkcie
\(g\colon y = |x + 1|- 2\)?}
{}{}{}{}{}{}}
\LANGes{
\Question{
Identifica la posible gráfica de la función
\(g(x) = |x + 1|- 2\).}
{}{}{}{}{}{}}

\IMGanswer{000040_05-figure2.pdf}
\IMGanswer{000040_05-figure0.pdf}
\IMGanswer{000040_05-figure1.pdf}
\IMGanswer{000040_05-figure3.pdf}\End

\Data\ID{9100004006}
\Updated{2020-11-24 02:11:20}
\Level{A}
\SubArea{Functions with absolute values}
\LANGen{
\Question{
Identify a possible graph of the function
\(g(x) = 2|x|- 1\).}
{}{}{}{}{}{}}
\LANGpl{
\Question{
Określ, który z rysunków przedstawia wykres funkcji
\(g\colon y = 2|x|- 1\).}
{}{}{}{}{}{}}
\LANGcs{
\Question{
Určete, který z obrázků je grafem funkce
\(g\colon y = 2|x|- 1\).}
{}{}{}{}{}{}}
\LANGsk{
\Question{
Ktorý z grafov je grafom funkcie
\(g\colon y = 2|x|- 1\)?}
{}{}{}{}{}{}}
\LANGes{
\Question{
Identifica la posible gráfica de la función
\(g(x) = 2|x|- 1\).}
{}{}{}{}{}{}}

\IMGanswer{000040_06-figure2.pdf}
\IMGanswer{000040_06-figure0.pdf}
\IMGanswer{000040_06-figure1.pdf}
\IMGanswer{000040_06-figure3.pdf}\End

\Data\ID{9100004007}
\Updated{2020-11-24 02:11:37}
\Level{A}
\SubArea{Functions with absolute values}
\LANGen{
\Question{
Identify a possible graph of the function
\(g(x) = -2|x - 1|\).}
{}{}{}{}{}{}}
\LANGpl{
\Question{
Określ, który z rysunków przedstawia wykres funkcji
\(g\colon y = -2|x - 1|\).}
{}{}{}{}{}{}}
\LANGcs{
\Question{
Určete, který z obrázků je grafem funkce
\(g\colon y = -2|x - 1|\).}
{}{}{}{}{}{}}
\LANGsk{
\Question{
Ktorý z grafov je grafom funkcie
\(g\colon y = -2|x - 1|\)?}
{}{}{}{}{}{}}
\LANGes{
\Question{
Identifica la posible gráfica de la función
\(g(x) = -2|x - 1|\).}
{}{}{}{}{}{}}

\IMGanswer{000040_07-figure0.pdf}
\IMGanswer{000040_07-figure1.pdf}
\IMGanswer{000040_07-figure2.pdf}
\IMGanswer{000040_07-figure3.pdf}\End

\Data\ID{9100004008}
\Updated{2020-11-24 02:11:06}
\Level{A}
\SubArea{Functions with absolute values}
\LANGen{
\Question{
Identify a possible graph of the function
\(g(x) = 2|x + 1|- 2\).}
{}{}{}{}{}{}}
\LANGpl{
\Question{
Określ, który z rysunków przedstawia wykres funkcji
\(g\colon y = 2|x + 1|- 2\).}
{}{}{}{}{}{}}
\LANGcs{
\Question{
Určete, který z obrázků je grafem funkce
\(g\colon y = 2|x + 1|- 2\).}
{}{}{}{}{}{}}
\LANGsk{
\Question{
Ktorý z grafov je grafom funkcie
\(g\colon y = 2|x + 1|- 2\)?}
{}{}{}{}{}{}}
\LANGes{
\Question{
Identifica la posible gráfica de la función
\(g(x) = 2|x + 1|- 2\).}
{}{}{}{}{}{}}

\IMGanswer{000040_08-figure1.pdf}
\IMGanswer{000040_08-figure0.pdf}
\IMGanswer{000040_08-figure2.pdf}
\IMGanswer{000040_08-figure3.pdf}\End

\Data\ID{9100024401}
\Updated{2022-05-29 12:05:45}
\Level{A}
\SubArea{Functions with absolute values}
\LANGen{
\Question{
Identify a possible graph of the function
\(f(x)= |x|-\sqrt{3}\).}
{}{}{}{}{}{}}
\LANGpl{
\Question{
Określ, który z poniższych rysunków przedstawia wykres funkcji
\(f\colon y = |x|-\sqrt{3}\).}
{}{}{}{}{}{}}
\LANGcs{
\Question{
Na kterém obrázku je znázorněn graf funkce
\(f\) dané
předpisem \(f\colon y = |x|-\sqrt{3}\)?}
{}{}{}{}{}{}}
\LANGsk{
\Question{
Ktorý z nasledujúcich grafov je grafom funkcie
\(f\colon y = |x|-\sqrt{3}\)?}
{}{}{}{}{}{}}
\LANGes{
\Question{
Identifica la posible gráfica de la función
\(f(x)= |x|-\sqrt{3}\).}
{}{}{}{}{}{}}

\IMGanswer{000244_01-figure3.pdf}
\IMGanswer{000244_01-figure0.pdf}
\IMGanswer{000244_01-figure1.pdf}
\IMGanswer{000244_01-figure2.pdf}\End

\Data\ID{9100024402}
\Updated{2020-11-24 02:11:52}
\Level{A}
\SubArea{Functions with absolute values}
\LANGen{
\Question{
Identify a possible graph of the function
\(f(x) = |x -\sqrt{3}|\).}
{}{}{}{}{}{}}
\LANGpl{
\Question{
Określ, który z poniższych rysunków przedstawia wykres funkcji
\(f\colon y = |x -\sqrt{3}|\).}
{}{}{}{}{}{}}
\LANGcs{
\Question{
Na kterém obrázku je znázorněn graf funkce
\(f\) dané
předpisem \(f\colon y = |x -\sqrt{3}|\)?}
{}{}{}{}{}{}}
\LANGsk{
\Question{
Ktorý z nasledujúcich grafov je grafom funkcie
\(f\colon y = |x -\sqrt{3}|\)?}
{}{}{}{}{}{}}
\LANGes{
\Question{
Identifica la posible gráfica de la función
\(f(x) = |x -\sqrt{3}|\).}
{}{}{}{}{}{}}

\IMGanswer{000244_02-figure0.pdf}
\IMGanswer{000244_02-figure1.pdf}
\IMGanswer{000244_02-figure2.pdf}
\IMGanswer{000244_02-figure3.pdf}\End

\Data\ID{9100024409}
\Updated{2020-11-24 02:11:03}
\Level{A}
\SubArea{Functions with absolute values}
\LANGen{
\Question{
Identify a possible graph of the function
\(f(x) = -|3 - x|- 1\).}
{}{}{}{}{}{}}
\LANGpl{
\Question{
Określ, który z rysunków przedstawia wykres funkcji
\(f\colon y = -|3 - x|- 1\)?}
{}{}{}{}{}{}}
\LANGcs{
\Question{
Je dána funkce \(f\colon y = -|3 - x|- 1\).
Vyberte její graf.}
{}{}{}{}{}{}}
\LANGsk{
\Question{
Ktorý z grafov je grafom funkcie
\(f\colon y = -|3 - x|- 1\).}
{}{}{}{}{}{}}
\LANGes{
\Question{
Identifica la posible gráfica de la función
\(f(x) = -|3 - x|- 1\).}
{}{}{}{}{}{}}

\IMGanswer{000244_09-figure2.pdf}
\IMGanswer{000244_09-figure0.pdf}
\IMGanswer{000244_09-figure1.pdf}
\IMGanswer{000244_09-figure3.pdf}\End

\Data\ID{9100024410}
\Updated{2020-11-24 02:11:17}
\Level{A}
\SubArea{Functions with absolute values}
\LANGen{
\Question{
Identify a possible graph of the function
\(f(x) = -|x - 4| + 2\).}
{}{}{}{}{}{}}
\LANGpl{
\Question{
Określ, który z rysunków przedstawia wykres funkcji
\(f\colon y = -|x - 4| + 2\)?}
{}{}{}{}{}{}}
\LANGcs{
\Question{
Je dána funkce \(f\colon y = -|x - 4| + 2\).
Vyberte její graf.}
{}{}{}{}{}{}}
\LANGsk{
\Question{
Ktorý z grafov je grafom funkcie
\(f\colon y = -|x - 4| + 2\).}
{}{}{}{}{}{}}
\LANGes{
\Question{
Identifica la posible gráfica de la función
\(f(x) = -|x - 4| + 2\).}
{}{}{}{}{}{}}

\IMGanswer{000244_10-figure0.pdf}
\IMGanswer{000244_10-figure1.pdf}
\IMGanswer{000244_10-figure2.pdf}
\IMGanswer{000244_10-figure3.pdf}\End

\Data\ID{9100024501}
\Updated{2020-11-24 02:11:52}
\Level{A}
\SubArea{Functions with absolute values}
\LANGen{
\Question{
Identify a possible graph of the function
\(f(x) = |x| + \sqrt{2}\).}
{}{}{}{}{}{}}
\LANGpl{
\Question{
Określ, który z poniższych rysunków przedstawia wykres funkcji
\(f\colon y = |x| + \sqrt{2}\).}
{}{}{}{}{}{}}
\LANGcs{
\Question{
Je dána funkce \(f\colon y = |x| + \sqrt{2}\).
Vyberte její graf.}
{}{}{}{}{}{}}
\LANGsk{
\Question{
Ktorý z nasledujúcich grafov je grafom funkcie
\(f\colon y = |x| + \sqrt{2}\)?}
{}{}{}{}{}{}}
\LANGes{
\Question{
Identifica la posible gráfica de la función
\(f(x) = |x| + \sqrt{2}\).}
{}{}{}{}{}{}}

\IMGanswer{000245_01-figure1.pdf}
\IMGanswer{000245_01-figure0.pdf}
\IMGanswer{000245_01-figure2.pdf}
\IMGanswer{000245_01-figure3.pdf}\End

\Data\ID{9100024502}
\Updated{2020-11-24 02:11:02}
\Level{A}
\SubArea{Functions with absolute values}
\LANGen{
\Question{
Identify a possible graph of the function
\(f(x) = |x + \sqrt{2}|\).}
{}{}{}{}{}{}}
\LANGpl{
\Question{
Określ, który z poniższych rysunków przedstawia wykres funkcji
\(f\colon y = |x + \sqrt{2}|\).}
{}{}{}{}{}{}}
\LANGcs{
\Question{
Je dána funkce \(f\colon y = |x + \sqrt{2}|\).
Vyberte její graf.}
{}{}{}{}{}{}}
\LANGsk{
\Question{
Ktorý z nasledujúcich grafov je grafom funkcie
\(f\colon y = |x + \sqrt{2}|\)?}
{}{}{}{}{}{}}
\LANGes{
\Question{
Identifica la posible gráfica de la función
\(f(x) = |x + \sqrt{2}|\).}
{}{}{}{}{}{}}

\IMGanswer{000245_02-figure3.pdf}
\IMGanswer{000245_02-figure0.pdf}
\IMGanswer{000245_02-figure1.pdf}
\IMGanswer{000245_02-figure2.pdf}\End

\Data\ID{9100024509}
\Updated{2020-11-24 02:11:12}
\Level{A}
\SubArea{Functions with absolute values}
\LANGen{
\Question{
Identify a possible graph of the function
\(f(x) = -|3 + x| + 1\).}
{}{}{}{}{}{}}
\LANGpl{
\Question{
Określ, który z rysunków przedstawia wykres funkcji
\(f\colon y = -|3 + x| + 1\).}
{}{}{}{}{}{}}
\LANGcs{
\Question{
Je dána funkce \(f\colon y = -|3 + x| + 1\).
Vyberte její graf.}
{}{}{}{}{}{}}
\LANGsk{
\Question{
Ktorý z grafov je grafom funkcie
\(f\colon y= -|3 + x| + 1\)?}
{}{}{}{}{}{}}
\LANGes{
\Question{
Identifica la posible gráfica de la función
\(f(x) = -|3 + x| + 1\).}
{}{}{}{}{}{}}

\IMGanswer{000245_09-figure3.pdf}
\IMGanswer{000245_09-figure0.pdf}
\IMGanswer{000245_09-figure1.pdf}
\IMGanswer{000245_09-figure2.pdf}\End

\Data\ID{9100024510}
\Updated{2020-11-24 02:11:23}
\Level{A}
\SubArea{Functions with absolute values}
\LANGen{
\Question{
Identify a possible graph of the function
\(f(x)= -|4 + x|- 2\).}
{}{}{}{}{}{}}
\LANGpl{
\Question{
Określ, który z rysunków przedstawia wykres funkcji
\(f\colon y = -|4 + x|- 2\).}
{}{}{}{}{}{}}
\LANGcs{
\Question{
Je dána funkce \(f\colon y = -|4 + x|- 2\).
Vyberte její graf.}
{}{}{}{}{}{}}
\LANGsk{
\Question{
Ktorý z grafov je grafom funkcie
\(f\colon y = -|4 + x|- 2\)?}
{}{}{}{}{}{}}
\LANGes{
\Question{
Identifica la posible gráfica de la función
\(f(x)= -|4 + x|- 2\).}
{}{}{}{}{}{}}

\IMGanswer{000245_10-figure1.pdf}
\IMGanswer{000245_10-figure0.pdf}
\IMGanswer{000245_10-figure2.pdf}
\IMGanswer{000245_10-figure3.pdf}\End

\Data\ID{1003019303}
\Updated{2022-08-25 10:08:58}
\Level{B}
\SubArea{Functions with absolute values}
\LANGen{
\Question{
Let \( f(x)=|2-x|+|x+1| \). Identify which of the statements is correct.}
{Function \( f \) has minimum at \( x=-1 \) and at \( x=2 \).}{Function \( f \) has minimum at \( x=-1 \) and maximum at \( x=2 \).}{Function \( f \) has maximum at \( x=-1 \) and at \( x=2 \).}{Function \( f \) has minimum at \( x=3 \).}{}{}}
\LANGpl{
\Question{
Niech \( f(x)=|2-x|+|x+1| \). Które z poniższych stwierdzeń jest prawdziwe?}
{Funkcja \( f \) osiąga minimum w  \( x=-1 \) i w \( x=2 \).}{Funkcja \( f \) osiąga minimum w  \( x=-1 \) i maksimum w \( x=2 \).}{Funkcja \( f \) osiąga maksimum w \( x=-1 \) i w  \( x=2 \).}{Funkcja \( f \) osiąga minimum w  \( x=3 \).}{}{}}
\LANGcs{
\Question{
Funkce \(f\) je dána předpisem \( f(x)=|2-x|+|x+1| \). Vyberte pravdivý výrok:}
{Funkce \( f \) má minimum v bodě \( -1 \) a v bodě \( 2 \).}{Funkce \( f \) má minimum v bodě \( -1 \) a maximum v bodě \( 2 \).}{Funkce \( f \) má maximum v bodě \( -1 \) a v bodě \( 2 \).}{Funkce \( f \) má minimum v bodě \( 3 \).}{}{}}
\LANGsk{
\Question{
Daná je funkcia \( f(x)=|2-x|+|x+1| \). Ktoré z nasledujúcich tvrdení je pravdivé?}
{Funkcia \( f \) má minimum v bode \( -1 \) a v bode \( 2 \).}{Funkcia \( f \) má minimum v bode \( -1 \) a maximum v bode \( 2 \).}{Funkcia \( f \) má maximum v bode \( -1 \) a v bode \( 2 \).}{Funkcia \( f \) má minimum v bode \( 3 \).}{}{}}
\LANGes{
\Question{
Sea \( f(x)=|2-x|+|x+1| \). Identifica cuál de las declaraciones es correcta.}
{La función \( f \) tiene un mínimo en \( x=-1 \) y en \( x=2 \).}{La función \( f \) tiene un mínimo en \( x=-1 \) y un máximo en \( x=2 \).}{La función \( f \) tiene un máximo en \( x=-1 \) y en \( x=2 \).}{La función \( f \) tiene un mínimo en \( x=3 \).}{}{}}
\End

\Data\ID{1003030905}
\Updated{2022-08-25 10:08:38}
\Level{B}
\SubArea{Functions with absolute values}
\LANGen{
\Question{
Let \( f(x)=|x-1|-2|x| \). Identify which of the following statements is true.}
{The function \( f \) is bounded above and is not bounded below.}{The function \( f \) is bounded below and is not bounded above.}{The function \( f \) is bounded.}{The function \( f \) is neither bounded above nor bounded below.}{}{}}
\LANGpl{
\Question{
Załóżmy, że  \( f(x)=|x-1|-2|x| \). Które z poniższych zdań jest prawdziwe?}
{Funkcja \( f \) jest ograniczona z góry i nie jest ograniczona z dołu.}{Funkcja \( f \) jest ograniczona z dołu i nie jest ograniczona z góry.}{Funkcja \( f \) jest ograniczona.}{Funkcja \( f \)  nie jest ograniczona ani z góry ani z dołu.}{}{}}
\LANGcs{
\Question{
Funkce \( f \) je dána předpisem \( f(x)=|x-1|-2|x| \). Vyberte pravdivý výrok.}
{Funkce  \( f \) je shora omezená a není zdola omezená.}{Funkce  \( f \) je zdola omezená a není shora omezená.}{Funkce  \( f \) je omezená.}{Funkce  \( f \) není shora ani zdola omezená.}{}{}}
\LANGsk{
\Question{
Daná je funkcia predpisom \( f(x)=|x-1|-2|x| \). Vyberte pravdivé tvrdenie o funkcii \( f \).}
{Funkcia \( f \) je ohraničená zhora a nie je ohraničená zdola.}{Funkcia \( f \) je ohraničená zdola a nie je ohraničená zhora.}{Funkcia \( f \) je ohraničená.}{Funkcia \( f \) nie je ohraničená zdola ani zhora.}{}{}}
\LANGes{
\Question{
Sea \( f(x)=|x-1|-2|x| \). Identifica cuál de las siguientes declaraciones es correcta.}
{La función \( f \) está acotada superiormente y no está acotada inferiormente.}{La función \( f \) está acotada inferiormente y no está acotada superiormente.}{La función \( f \) está acotada.}{La función \( f \) no está acotada inferiormente ni superiormente.}{}{}}
\End

\Data\ID{1003049201}
\Updated{2020-11-24 08:11:33}
\Level{B}
\SubArea{Functions with absolute values}
\LANGen{
\Question{
Let \( f(x)=2|2x+3|-|x|\left|\frac12-x\right|\). Identify which of the following function values is the highest.}
{\( f(2) \)}{\( f(0) \)}{\( f(-1.5) \)}{\( f(3.5) \)}{}{}}
\LANGpl{
\Question{
Niech  \( f(x)=2|2x+3|-|x|\left|\frac12-x\right|\). Określ, która z podanych wartości funkcji jest najwyższa.}
{\( f(2) \)}{\( f(0) \)}{\( f(-1{,}5) \)}{\( f(3{,}5) \)}{}{}}
\LANGcs{
\Question{
Funkce f je dána předpisem \( f(x)=2|2x+3|-|x|\left|\frac12-x\right|\). Která z daných funkčních hodnot je největší?}
{\( f(2) \)}{\( f(0) \)}{\( f(-1{,}5) \)}{\( f(3{,}5) \)}{}{}}
\LANGsk{
\Question{
Daná je funkcia \( f(x)=2|2x+3|-|x|\left|\frac12-x\right|\). Ktorá z nasledujúcich funkčných hodnôt je najväčšia?}
{\( f(2) \)}{\( f(0) \)}{\( f(-1{,}5) \)}{\( f(3{,}5) \)}{}{}}
\LANGes{
\Question{
Sea \( f(x)=2|2x+3|-|x|\left|\frac12-x\right|\). Identifica cuál de los valores funcionales es el más grande.}
{\( f(2) \)}{\( f(0) \)}{\( f(-1.5) \)}{\( f(3.5) \)}{}{}}
\End

\Data\ID{1003049301}
\Updated{2021-04-11 08:04:09}
\Level{B}
\SubArea{Functions with absolute values}
\LANGen{
\Question{
Let \( f(x)=\frac12 (|x|-3) \). Identify which of the following statements is false.}
{The range of the function \( f \) is the interval \( [-3;\infty ) \).}{The domain of the function \( f \) is the interval \( (-\infty;\infty) \).}{The function \( f \) is even.}{The function \( f \) is bounded below.}{}{}}
\LANGpl{
\Question{
Załóżmy, że \( f(x)=\frac12 (|x|-3) \). Określ, które z poniższych zdań jest fałszywe.}
{Zakres funkcji \( f \) mieści się w przedziale \( \langle-3;\infty ) \).}{Dziedziną funkcji \( f \) jest przedział \( (-\infty;\infty) \).}{Funkcja  \( f \) jest parzysta.}{Funkcja \( f \) jest ograniczona z dołu.}{}{}}
\LANGcs{
\Question{
Funkce \( f \)  je dána předpisem \( f(x)=\frac12 (|x|-3) \). Vyberte nepravdivý výrok.}
{Oborem hodnot funkce \( f \) je interval \( \langle-3;\infty ) \).}{Definičním oborem funkce \( f \) je interval \( (-\infty;\infty) \).}{Funkce \( f \) je sudá.}{Funkce \( f \) je zdola omezená.}{}{}}
\LANGsk{
\Question{
Daná je funkcia predpisom \( f(x)=\frac12 (|x|-3) \). Vyberte nepravdivé tvrdenie o funkcii \( f \).}
{Obor hodnôt funkcie \( f \) je interval \( \langle-3;\infty ) \).}{Definičný obor funkcie \( f \) je interval \( (-\infty;\infty) \).}{Funkcia \( f \) je párna.}{Funkcia \( f \) je ohraničená zdola.}{}{}}
\LANGes{
\Question{
Sea \( f(x)=\frac12 (|x|-3) \). Identifica cuál de las declaraciones es falsa.}
{El Rango de la función \( f \) es el intervalo \( [-3;\infty ) \).}{El Dominio de la función \( f \) es el intervalo \( (-\infty;\infty) \).}{La función \( f \) es par.}{La función \( f \) es acotada inferiormente.}{}{}}
\End

\Data\ID{1003049302}
\Updated{2020-11-24 08:11:58}
\Level{B}
\SubArea{Functions with absolute values}
\LANGen{
\Question{
Let \( f(x)=2|-x|+1 \). Identify which of the following statements is false.}
{The function \( f \) is an injective (one-to-one) function.}{The function \( f \) is increasing in the interval \( [0; \infty) \).}{The function \( f \) has minimum at \( x=0 \).}{The function \( f \) is non-increasing in the interval \( (-\infty; 0] \).}{}{}}
\LANGpl{
\Question{
Załóżmy, że \( f(x)=2|-x|+1 \). Które z poniższych zdań jest fałszywe?}
{Funkcja \( f \) jest iniektywna.}{Funkcja \( f \) jest rosnąca w przedziale \( \langle0; \infty) \).}{Funkcja \( f \) osiąga minimum w \( x=0 \).}{Funkcja \( f \) jest nierosnąca w przedziale \( (-\infty; 0\rangle \).}{}{}}
\LANGcs{
\Question{
Funkce \( f \) je dána předpisem \( f(x)=2|-x|+1 \). Vyberte nepravdivý výrok.}
{Funkce \( f \) je prostá.}{Funkce \( f \) je rostoucí v intervalu \( \langle0; \infty) \).}{Funkce \( f \) má minimum v bodě \( x=0 \).}{Funkce \( f \) je nerostoucí v intervalu \( (-\infty; 0\rangle \).}{}{}}
\LANGsk{
\Question{
Daná je funkcia predpisom \( f(x)=2|-x|+1 \). Vyberte nepravdivé tvrdenie o funkcii \( f \).}
{Funkcia \( f \) je prostá.}{Funkcia \( f \) je rastúca na intervale \( \langle0; \infty) \).}{Funkcia \( f \) má minimum v bode \( x=0 \).}{Funkcia \( f \) je nerastúca na intervale \( (-\infty; 0\rangle \).}{}{}}
\LANGes{
\Question{
Sea \( f(x)=2|-x|+1 \). Identifica cuál de las siguientes declaraciones es falsa.}
{La función \( f \) es inyectiva.}{La función \( f \) es creciente en el intervalo \( [0; \infty) \).}{La función \( f \) tiene mínimo en \( x=0 \).}{La función \( f \) es no creciente en el intervalo \( (-\infty; 0] \).}{}{}}
\End

\Data\ID{1003049303}
\Updated{2021-04-11 08:04:07}
\Level{B}
\SubArea{Functions with absolute values}
\LANGen{
\Question{
Let \( f(x)=-|x+3| \). Identify which of the following statements is false.}
{The function \( f \) is even.}{The function \( f \) is decreasing in the interval \( [3; \infty) \).}{The function \( f \) has maximum at \( x=-3 \).}{The range of the function \( f \) is the interval \( (-\infty; 0] \).}{}{}}
\LANGpl{
\Question{
Załóżmy, że  \( f(x)=-|x+3| \). Które z poniższych stwierdzeń jest fałszywe?}
{Funkcja  \( f \) jest parzysta.}{Funkcja \( f \) jest malejąca w przedziale \( \langle3; \infty) \).}{Funkcja  \( f \) osiąga maksimum w \( x=-3 \).}{Zakresem funkcji  \( f \) jest przedział \( (-\infty; 0\rangle \).}{}{}}
\LANGcs{
\Question{
Funkce \( f \) je dána předpisem \( f(x)=-|x+3| \). Vyberte nepravdivý výrok.}
{Funkce \( f \) je sudá.}{Funkce \( f \) je klesající v intervalu \( \langle3; \infty) \).}{Funkce \( f \) má maximum v \( x=-3 \).}{Oborem hodnot funkce \( f \) je interval \( (-\infty; 0\rangle \).}{}{}}
\LANGsk{
\Question{
Daná je funkcia predpisom \( f(x)=-|x+3| \). Vyberte nepravdivé tvrdenie o funkcii \( f \).}
{Funkcia \( f \) je párna.}{Funkcia \( f \) je klesajúca na intervale \( \langle3; \infty) \).}{Funkcia \( f \) má maximum v bode \( x=-3 \).}{Obor hodnôt funkcie \( f \) je interval \( (-\infty; 0\rangle \).}{}{}}
\LANGes{
\Question{
Sea \( f(x)=-|x+3| \). Identifica cuál de las siguientes declaraciones es falsa.}
{La función \( f \) es par.}{La función \( f \) es decreciente en el intervalo \( [3; \infty) \).}{La función \( f \) tiene máximo en \( x=-3 \).}{El Rango de la función \( f \) es el intervalo \( (-\infty; 0] \).}{}{}}
\End

\Data\ID{1003049304}
\Updated{2020-11-24 09:11:45}
\Level{B}
\SubArea{Functions with absolute values}
\LANGen{
\Question{
Identify which of the following functions is increasing in the  interval \( (-\infty; +2] \).}
{\( m(x)=-2|5-x|+1 \)}{\( h(x)=|x|+2 \)}{\( g(x)=-|x|-2 \)}{\( f(x)=-|x+2| \)}{}{}}
\LANGpl{
\Question{
Która z poniższych funkcji jest rosnąca w przedziale \( (-\infty; +2\rangle \)?}
{\( m(x)=-2|5-x|+1 \)}{\( h(x)=|x|+2 \)}{\( g(x)=-|x|-2 \)}{\( f(x)=-|x+2| \)}{}{}}
\LANGcs{
\Question{
Která z následujících funkcí je rostoucí v intervalu \( (-\infty; +2\rangle \)?}
{\( m(x)=-2|5-x|+1 \)}{\( h(x)=|x|+2 \)}{\( g(x)=-|x|-2 \)}{\( f(x)=-|x+2| \)}{}{}}
\LANGsk{
\Question{
Ktorá z nasledujúcich funkcií je rastúca na intervale \( (-\infty; +2\rangle \)?}
{\( m(x)=-2|5-x|+1 \)}{\( h(x)=|x|+2 \)}{\( g(x)=-|x|-2 \)}{\( f(x)=-|x+2| \)}{}{}}
\LANGes{
\Question{
Identifica cuál de las siguientes funciones es creciente en el intervalo \( (-\infty; +2] \).}
{\( m(x)=-2|5-x|+1 \)}{\( h(x)=|x|+2 \)}{\( g(x)=-|x|-2 \)}{\( f(x)=-|x+2| \)}{}{}}
\End

\Data\ID{1003049305}
\Updated{2021-04-11 08:04:33}
\Level{B}
\SubArea{Functions with absolute values}
\LANGen{
\Question{
Identify a real number \( a \) such  that the function \( f(x)=a|-12-ax|-a \) has a maximum at \( x=3 \).}
{\( -4 \)}{\( 3 \)}{\( 4 \)}{There is no such number.}{}{}}
\LANGpl{
\Question{
Wyznacz liczbę rzeczywistą \( a \) taką, dla której funkcja \( f(x)=a|-12-ax|-a \) osiąga maksimum w \( x=3 \).}
{\( -4 \)}{\( 3 \)}{\( 4 \)}{Taka liczba nie istnieje.}{}{}}
\LANGcs{
\Question{
Vyberte reálné číslo \( a \) tak, aby funkce \( f \), která je dána předpisem  \( f(x)=a|-12-ax|-a \)  měla maximum v bodě \( x=3 \).}
{\( -4 \)}{\( 3 \)}{\( 4 \)}{Takové číslo neexistuje.}{}{}}
\LANGsk{
\Question{
Vyberte také reálne číslo \( a \), aby funkcia daná predpisom \( f(x)=a|-12-ax|-a \) mala maximum v bode \( x=3 \).}
{\( -4 \)}{\( 3 \)}{\( 4 \)}{Také číslo \( a \) neexistuje.}{}{}}
\LANGes{
\Question{
Identifica el número real \( a \) tal que la función \( f(x)=a|-12-ax|-a \) tenga un máximo en \( x=3 \).}
{\( -4 \)}{\( 3 \)}{\( 4 \)}{No existe tal número.}{}{}}
\End

\Data\ID{1003049306}
\Updated{2020-11-24 09:11:25}
\Level{B}
\SubArea{Functions with absolute values}
\LANGen{
\Question{
Identify which of the following functions is odd.}
{\( g(x)=|x|-|-x| \)}{\( h(x)=|x|-x \)}{\( f(x)=-|-x| \)}{\( m(x)=|-1| \)}{}{}}
\LANGpl{
\Question{
Określ, która z poniższych funkcji jest nieparzysta.}
{\( g(x)=|x|-|-x| \)}{\( h(x)=|x|-x \)}{\( f(x)=-|-x| \)}{\( m(x)=|-1| \)}{}{}}
\LANGcs{
\Question{
Která z následujících funkcí je lichá?}
{\( g(x)=|x|-|-x| \)}{\( h(x)=|x|-x \)}{\( f(x)=-|-x| \)}{\( m(x)=|-1| \)}{}{}}
\LANGsk{
\Question{
Ktorá z nasledujúcich funkcií je nepárna?}
{\( g(x)=|x|-|-x| \)}{\( h(x)=|x|-x \)}{\( f(x)=-|-x| \)}{\( m(x)=|-1| \)}{}{}}
\LANGes{
\Question{
¿Cuál de las funciones es impar?}
{\( g(x)=|x|-|-x| \)}{\( h(x)=|x|-x \)}{\( f(x)=-|-x| \)}{\( m(x)=|-1| \)}{}{}}
\End

\Data\ID{1003072701}
\Updated{2020-11-24 09:11:22}
\Level{B}
\SubArea{Functions with absolute values}
\LANGen{
\Question{
Which of the following functions is neither even nor odd?}
{\( h(x)=|x+1|+|x+1| \)}{\( f(x)=|x+1|+|x-1| \)}{\( g(x)=|x+1|-|x-1| \)}{\( m(x)=|x-1|-|x+1| \)}{}{}}
\LANGpl{
\Question{
Która z podanych funkcji nie jest ani parzysta ani nieparzysta?}
{\( h(x)=|x+1|+|x+1| \)}{\( f(x)=|x+1|+|x-1| \)}{\( g(x)=|x+1|-|x-1| \)}{\( m(x)=|x-1|-|x+1| \)}{}{}}
\LANGcs{
\Question{
Která z následujících funkcí není ani sudá ani lichá?}
{\( h(x)=|x+1|+|x+1| \)}{\( f(x)=|x+1|+|x-1| \)}{\( g(x)=|x+1|-|x-1| \)}{\( m(x)=|x-1|-|x+1| \)}{}{}}
\LANGsk{
\Question{
Ktorá z následujúcich funkcií nie je párna ani nepárna?}
{\( h(x)=|x+1|+|x+1| \)}{\( f(x)=|x+1|+|x-1| \)}{\( g(x)=|x+1|-|x-1| \)}{\( m(x)=|x-1|-|x+1| \)}{}{}}
\LANGes{
\Question{
¿Cuál de las funciones no es par ni impar?}
{\( h(x)=|x+1|+|x+1| \)}{\( f(x)=|x+1|+|x-1| \)}{\( g(x)=|x+1|-|x-1| \)}{\( m(x)=|x-1|-|x+1| \)}{}{}}
\End

\Data\ID{1003072702}
\Updated{2022-04-23 05:04:37}
\Level{B}
\SubArea{Functions with absolute values}
\LANGen{
\Question{
Which of the following functions is bounded?}
{\( g(x)=|x-2|-|x| \)}{\( f(x)=|x-2|+|x| \)}{\( h(x)=|x+2|+|x| \)}{\( m(x)=|2x|-|x| \)}{}{}}
\LANGpl{
\Question{
Która z poniższych funkcji jest ograniczona?}
{\( g(x)=|x-2|-|x| \)}{\( f(x)=|x-2|+|x| \)}{\( h(x)=|x+2|+|x| \)}{\( m(x)=|2x|-|x| \)}{}{}}
\LANGcs{
\Question{
Která z následujících funkcí je omezená?}
{\( g(x)=|x-2|-|x| \)}{\( f(x)=|x-2|+|x| \)}{\( h(x)=|x+2|+|x| \)}{\( m(x)=|2x|-|x| \)}{}{}}
\LANGsk{
\Question{
Ktorá z následujúcich funkcií je ohraničená?}
{\( g(x)=|x-2|-|x| \)}{\( f(x)=|x-2|+|x| \)}{\( h(x)=|x+2|+|x| \)}{\( m(x)=|2x|-|x| \)}{}{}}
\LANGes{
\Question{
¿Cuál de las funciones es acotada?}
{\( g(x)=|x-2|-|x| \)}{\( f(x)=|x-2|+|x| \)}{\( h(x)=|x+2|+|x| \)}{\( m(x)=|2x|-|x| \)}{}{}}
\End

\Data\ID{1003072703}
\Updated{2020-11-24 09:11:54}
\Level{B}
\SubArea{Functions with absolute values}
\LANGen{
\Question{
Find the true statement about the domain or the range of the function \( f(x)=|x+2|-|1-x| \).}
{\( H(f)=[-3;3]  \)}{\( D(f)=[-3;3] \)}{\( H(f)=\mathbb{R}  \)}{\( D(f)=[-2;1] \)}{}{}}
\LANGpl{
\Question{
Wybierz prawdziwe stwierdzenie dotyczące dziedziny lub zakresu funkcji  \( f(x)=|x+2|-|1-x| \).}
{\( H(f)=\langle-3;3\rangle  \)}{\( D(f)=\langle-3;3\rangle \)}{\( H(f)=\mathbb{R}  \)}{\( D(f)=\langle-2;1\rangle \)}{}{}}
\LANGcs{
\Question{
Funkce f je dána předpisem \( f(x)=|x+2|-|1-x| \). Vyberte pravdivý výrok o definičním oboru \(D(f)\), resp. oboru hodnot \(H(f)\) funkce f.}
{\( H(f)=\langle-3;3\rangle  \)}{\( D(f)=\langle-3;3\rangle \)}{\( H(f)=\mathbb{R}  \)}{\( D(f)=\langle-2;1\rangle \)}{}{}}
\LANGsk{
\Question{
Vyberte pravdivé tvrdenie o definičnom obore alebo obore hodnôt funkcie \( f(x)=|x+2|-|1-x| \).}
{\( H(f)=\langle-3;3\rangle  \)}{\( D(f)=\langle-3;3\rangle \)}{\( H(f)=\mathbb{R}  \)}{\( D(f)=\langle-2;1\rangle \)}{}{}}
\LANGes{
\Question{
¿Cuál de las declaraciones sobre el Dominio o el Rango de la función \( f(x)=|x+2|-|1-x| \) es verdadera?}
{\( H(f)=[-3;3]  \)}{\( D(f)=[-3;3] \)}{\( H(f)=\mathbb{R}  \)}{\( D(f)=[-2;1] \)}{}{}}
\End

\Data\ID{1003072704}
\Updated{2020-11-24 09:11:09}
\Level{B}
\SubArea{Functions with absolute values}
\LANGen{
\Question{
Find the true statement about the function \( f(x)=2|x-2|-|x-1|-x \).}
{The function \( f \) has minimum at \( x=3 \).}{The function \( f \) has maximum at \( x=1 \).}{The function \( f \) has minimum at \( x=-3 \).}{The function \( f \) has no minimum.}{}{}}
\LANGpl{
\Question{
Wybierz prawdziwe stwierdzenie dotyczące  funkcji \( f(x)=2|x-2|-|x-1|-x \).}
{Funkcja \( f \) ma minimum w  \( x=3 \).}{Funkcja \( f \) ma maksimum w \( x=1 \).}{Funkcja \( f \) ma minimum w  \( x=-3 \).}{Funkcja \( f \) nie ma minimum.}{}{}}
\LANGcs{
\Question{
Funkce f je dána předpisem \( f(x)=2|x-2|-|x-1|-x \). Vyberte pravdivý výrok.}
{Funkce \( f \) má minimum v bodě \( x=3 \).}{Funkce \( f \) má maximum v bodě \( x=1 \).}{Funkce \( f \) má minimum v bodě \( x=-3 \).}{Funkce \( f \) nemá minimum.}{}{}}
\LANGsk{
\Question{
Ktoré z tvrdení o funkcii \( f(x)=2|x-2|-|x-1|-x \) je pravdivé?}
{Funkcia \( f \) má minimum v \( x=3 \).}{Funkcia \( f \) má maximum v \( x=1 \).}{Funkcia \( f \) má minimum v  \( x=-3 \).}{Funkcia \( f \) nemá minimum.}{}{}}
\LANGes{
\Question{
¿Cuál de las declaraciones sobre la función \( f(x)=2|x-2|-|x-1|-x \) es verdadera?}
{La función \( f \) tiene mínimo en \( x=3 \).}{La función \( f \) tiene máximo \( x=1 \).}{La función \( f \) tiene mínimo en \( x=-3 \).}{La función \( f \) no tiene mínimo.}{}{}}
\End

\Data\ID{1003072705}
\Updated{2021-04-11 09:04:54}
\Level{B}
\SubArea{Functions with absolute values}
\LANGen{
\Question{
Find the false statement about the function \( f(x)=|x+1|-|x-2|+|x| \).}
{The function \( f \) is injective (one-to-one).}{The function \( f \) has minimum at \( x=-1 \).}{The function \( f \) is bounded below.}{The range of the function \( f \) is \( [-2;\infty) \).}{}{}}
\LANGpl{
\Question{
Wybierz fałszywe stwierdzenie dotyczące funkcji \( f(x)=|x+1|-|x-2|+|x| \).}
{Funkcja \( f \) jest iniektywna  (jeden-do-jednego).}{Funkcja \( f \) ma minimum w \( x=-1 \).}{Funkcja \( f \) jest ograniczona z dołu.}{Zakres funkcji \( f \) wynosi \( \langle-2;\infty) \).}{}{}}
\LANGcs{
\Question{
Funkce f je dána předpisem \( f(x)=|x+1|-|x-2|+|x| \). Vyberte nepravdivý výrok.}
{Funkce \( f \) je prostá.}{Funkce \( f \) má minimum v bodě \( x=-1 \).}{Funkce \( f \) je zdola omezená.}{Obor hodnot funkce \( f \) je \(H(f) = \langle-2;\infty) \).}{}{}}
\LANGsk{
\Question{
Ktoré z tvrdení o funkcii \( f(x)=|x+1|-|x-2|+|x| \) nie je pravdivé?}
{Funkcia \( f \) je prostá.}{Funkcia \( f \) má minimum v bode \( x=-1 \).}{Funkcia \( f \) je zdola ohraničená.}{Obor hodnôt funkcie \( f \) je \(H(f) = \langle-2;\infty) \).}{}{}}
\LANGes{
\Question{
Encuentra la declaración falsa sobre la función \( f(x)=|x+1|-|x-2|+|x| \).}
{La función \( f \) es inyectiva.}{La función \( f \) tiene un mínimo en \( x=-1 \).}{La función \( f \) es acotada inferiormente.}{El Rango de la función \( f \) es \( [-2;\infty) \).}{}{}}
\End

\Data\ID{1103072501}
\Updated{2022-04-23 05:04:37}
\Level{B}
\SubArea{Functions with absolute values}
\IMG{1103072501.pdf}
\LANGen{
\Question{
Function \( f \) is given by the graph. Identify which of the following statements is true.}
{\( f(x)=x-|x|;\ x\in[-4;4] \)}{\( f(x)=x+|x|;\ x\in[-4;4] \)}{\( f(x)=|x|-x;\ x\in[-4;4] \)}{\( f(x)=-x-|x|;\ x\in[-4;4] \)}{}{}}
\LANGpl{
\Question{
Funkcja \( f \) przedstawiona jest za pomocą wykresu. Które z poniższych stwierdzeń jest prawdziwe?}
{\( f(x)=x-|x|;\ x\in\langle-4;4\rangle \)}{\( f(x)=x+|x|;\ x\in\langle-4;4\rangle \)}{\( f(x)=|x|-x;\ x\in\langle-4;4\rangle \)}{\( f(x)=-x-|x|;\ x\in\langle-4;4\rangle \)}{}{}}
\LANGcs{
\Question{
Funkce \( f \) je dána grafem. Vyberte pravdivý výrok.}
{\( f(x)=x-|x|;\ x\in\langle-4;4\rangle \)}{\( f(x)=x+|x|;\ x\in\langle-4;4\rangle \)}{\( f(x)=|x|-x;\ x\in\langle-4;4\rangle \)}{\( f(x)=-x-|x|;\ x\in\langle-4;4\rangle \)}{}{}}
\LANGsk{
\Question{
Na obrázku je graf funkcie \( f \). Ktoré z nasledujúcich tvrdení je pravdivé?}
{\( f(x)=x-|x|;\ x\in\langle-4;4\rangle \)}{\( f(x)=x+|x|;\ x\in\langle-4;4\rangle \)}{\( f(x)=|x|-x;\ x\in\langle-4;4\rangle \)}{\( f(x)=-x-|x|;\ x\in\langle-4;4\rangle \)}{}{}}
\LANGes{
\Question{
La función \( f \) viene dada por la gráfica, identifica cuál de las declaraciones es correcta.}
{\( f(x)=x-|x|;\ x\in[-4;4] \)}{\( f(x)=x+|x|;\ x\in[-4;4] \)}{\( f(x)=|x|-x;\ x\in[-4;4] \)}{\( f(x)=-x-|x|;\ x\in[-4;4] \)}{}{}}
\End

\Data\ID{1103072502}
\Updated{2021-04-11 09:04:40}
\Level{B}
\SubArea{Functions with absolute values}
\IMG{1103072502.pdf}
\LANGen{
\Question{
Function \( f \) is given by the graph. Identify which of the following statements is true.}
{\( f(x)=|x+1|-x;\ x\in[-4;4] \)}{\( f(x)=|x|-x+1;\ x\in[-4;4] \)}{\( f(x)=x-|x+1|;\ x\in[-4;4] \)}{\( f(x)=x-|x|+1;\ x\in[-4;4] \)}{}{}}
\LANGpl{
\Question{
Funkcja \( f \) przedstawiona jest za pomocą wykresu. Które z poniższych stwierdzeń jest prawdziwe?}
{\( f(x)=|x+1|-x;\ x\in\langle-4;4\rangle \)}{\( f(x)=|x|-x+1;\ x\in\langle-4;4\rangle \)}{\( f(x)=x-|x+1|;\ x\in\langle-4;4\rangle \)}{\( f(x)=x-|x|+1;\ x\in\langle-4;4\rangle \)}{}{}}
\LANGcs{
\Question{
Funkce \( f \) je dána grafem. Vyberte pravdivý výrok.}
{\( f(x)=|x+1|-x;\ x\in\langle-4;4\rangle \)}{\( f(x)=|x|-x+1;\ x\in\langle-4;4\rangle \)}{\( f(x)=x-|x+1|;\ x\in\langle-4;4\rangle \)}{\( f(x)=x-|x|+1;\ x\in\langle-4;4\rangle \)}{}{}}
\LANGsk{
\Question{
Na obrázku je graf funkcie \( f \). Ktoré z nasledujúcich tvrdení je pravdivé?}
{\( f(x)=|x+1|-x;\ x\in\langle-4;4\rangle \)}{\( f(x)=|x|-x+1;\ x\in\langle-4;4\rangle \)}{\( f(x)=x-|x+1|;\ x\in\langle-4;4\rangle \)}{\( f(x)=x-|x|+1;\ x\in\langle-4;4\rangle \)}{}{}}
\LANGes{
\Question{
La función \( f \) viene dada por la gráfica. Identifica cuál de las declaraciones es correcta.}
{\( f(x)=|x+1|-x;\ x\in[-4;4] \)}{\( f(x)=|x|-x+1;\ x\in[-4;4] \)}{\( f(x)=x-|x+1|;\ x\in[-4;4] \)}{\( f(x)=x-|x|+1;\ x\in[-4;4] \)}{}{}}
\End

\Data\ID{1103072503}
\Updated{2022-04-23 05:04:37}
\Level{B}
\SubArea{Functions with absolute values}
\IMG{1103072503.pdf}
\LANGen{
\Question{
Function \( f \) is given by the graph. Identify which of the following statements is true.}
{\( f(x)=|x|-|x-2|;\ x\in[-4;4] \)}{\( f(x)=|x|-|x+2|;\ x\in[-4;4] \)}{\( f(x)=|x+2|-|x|;\ x\in[-4;4] \)}{\( f(x)=|x-2|-|x|;\ x\in[-4;4] \)}{}{}}
\LANGpl{
\Question{
Funkcja \( f \) przedstawiona jest za pomocą wykresu. Które z poniższych stwierdzeń jest prawdziwe?}
{\( f(x)=|x|-|x-2|;\ x\in\langle-4;4\rangle \)}{\( f(x)=|x|-|x+2|;\ x\in\langle-4;4\rangle \)}{\( f(x)=|x+2|-|x|;\ x\in\langle-4;4\rangle \)}{\( f(x)=|x-2|-|x|;\ x\in\langle-4;4\rangle \)}{}{}}
\LANGcs{
\Question{
Funkce \( f \) je dána grafem. Vyberte pravdivý výrok.}
{\( f(x)=|x|-|x-2|;\ x\in\langle-4;4\rangle \)}{\( f(x)=|x|-|x+2|;\ x\in\langle-4;4\rangle \)}{\( f(x)=|x+2|-|x|;\ x\in\langle-4;4\rangle \)}{\( f(x)=|x-2|-|x|;\ x\in\langle-4;4\rangle \)}{}{}}
\LANGsk{
\Question{
Na obrázku je graf funkcie \( f \). Ktoré z nasledujúcich tvrdení je pravdivé?}
{\( f(x)=|x|-|x-2|;\ x\in\langle-4;4\rangle \)}{\( f(x)=|x|-|x+2|;\ x\in\langle-4;4\rangle \)}{\( f(x)=|x+2|-|x|;\ x\in\langle-4;4\rangle \)}{\( f(x)=|x-2|-|x|;\ x\in\langle-4;4\rangle \)}{}{}}
\LANGes{
\Question{
La función \( f \) viene dada por la gráfica. Identifica cuál de las declaraciones es correcta.}
{\( f(x)=|x|-|x-2|;\ x\in[-4;4] \)}{\( f(x)=|x|-|x+2|;\ x\in[-4;4] \)}{\( f(x)=|x+2|-|x|;\ x\in[-4;4] \)}{\( f(x)=|x-2|-|x|;\ x\in[-4;4] \)}{}{}}
\End

\Data\ID{1103072504}
\Updated{2021-04-11 09:04:15}
\Level{B}
\SubArea{Functions with absolute values}
\IMG{1103072504.pdf}
\LANGen{
\Question{
Function \( f \) is given by the graph. Identify which of the following statements is false.}
{\( f(x)=|x-1|-|2x|;\ x\in[-4;4] \)}{\( f(x)=|2x|-|x-1|;\ x\in[-4;4] \)}{\( f(x)=|2x|-|1-x|;\ x\in[-4;4] \)}{\( f(x)=2|x|-|x-1|;\ x\in[-4;4] \)}{}{}}
\LANGpl{
\Question{
Funkcja \( f \) przedstawiona jest za pomocą wykresu. Które z poniższych stwierdzeń jest fałszywe?}
{\( f(x)=|x-1|-|2x|;\ x\in\langle-4;4\rangle \)}{\( f(x)=|2x|-|x-1|;\ x\in\langle-4;4\rangle \)}{\( f(x)=|2x|-|1-x|;\ x\in\langle-4;4\rangle \)}{\( f(x)=2|x|-|x-1|;\ x\in\langle-4;4\rangle \)}{}{}}
\LANGcs{
\Question{
Funkce \( f \) je dána grafem. Vyberte nepravdivý výrok.}
{\( f(x)=|x-1|-|2x|;\ x\in\langle-4;4\rangle \)}{\( f(x)=|2x|-|x-1|;\ x\in\langle-4;4\rangle \)}{\( f(x)=|2x|-|1-x|;\ x\in\langle-4;4\rangle \)}{\( f(x)=2|x|-|x-1|;\ x\in\langle-4;4\rangle \)}{}{}}
\LANGsk{
\Question{
Na obrázku je graf funkcie \( f \). Ktoré z nasledujúcich tvrdení je nepravdivé?}
{\( f(x)=|x-1|-|2x|;\ x\in\langle-4;4\rangle \)}{\( f(x)=|2x|-|x-1|;\ x\in\langle-4;4\rangle \)}{\( f(x)=|2x|-|1-x|;\ x\in\langle-4;4\rangle \)}{\( f(x)=2|x|-|x-1|;\ x\in\langle-4;4\rangle \)}{}{}}
\LANGes{
\Question{
La función \( f \) viene dada por la gráfica. Identifica cuál de las declaraciones es falsa.}
{\( f(x)=|x-1|-|2x|;\ x\in[-4;4] \)}{\( f(x)=|2x|-|x-1|;\ x\in[-4;4] \)}{\( f(x)=|2x|-|1-x|;\ x\in[-4;4] \)}{\( f(x)=2|x|-|x-1|;\ x\in[-4;4] \)}{}{}}
\End

\Data\ID{1103072505}
\Updated{2020-11-24 11:11:10}
\Level{B}
\SubArea{Functions with absolute values}
\LANGen{
\Question{
Identify which of the graphs represents the  function \( f(x)=|2x-4|-|x-1|;\ x\in [-4;4] \).}
{}{}{}{}{}{}}
\LANGpl{
\Question{
Z podanych wykresów wybierz ten, który przedstawia funkcję  \( f(x)=|2x-4|-|x-1|;\ x\in\langle-4;4\rangle \).}
{}{}{}{}{}{}}
\LANGcs{
\Question{
Na kterém obrázku je graf funkce \( f(x)=|2x-4|-|x-1|;\ x\in\langle-4;4\rangle \)?}
{}{}{}{}{}{}}
\LANGsk{
\Question{
Ktorý z grafov je graf funkcie \( f(x)=|2x-4|-|x-1|;\ x\in\langle-4;4\rangle \)?}
{}{}{}{}{}{}}
\LANGes{
\Question{
Identifica cuál de las gráficas representa la función \( f(x)=|2x-4|-|x-1|;\ x\in [-4;4] \).}
{}{}{}{}{}{}}

\IMGanswer{1103072505a.pdf}
\IMGanswer{1103072505b.pdf}
\IMGanswer{1103072505c.pdf}
\IMGanswer{1103072505d.pdf}\End

\Data\ID{1103072506}
\Updated{2020-11-24 11:11:08}
\Level{B}
\SubArea{Functions with absolute values}
\LANGen{
\Question{
Identify which of the graphs below represents  the function \( f(x)=x+1-|x+1|;\ x\in [-4;4] \).}
{}{}{}{}{}{}}
\LANGpl{
\Question{
Określ, który z poniższych wykresów przedstawia funkcję  \( f(x)=x+1-|x+1|;\ x\in\langle-4;4\rangle \).}
{}{}{}{}{}{}}
\LANGcs{
\Question{
Na kterém obrázku je graf funkce \( f(x)=x+1-|x+1|;\ x\in\langle-4;4\rangle \)?}
{}{}{}{}{}{}}
\LANGsk{
\Question{
Ktorý z nasledujúcich grafov je grafom funkcie \( f(x)=x+1-|x+1|;\ x\in\langle-4;4\rangle \)?}
{}{}{}{}{}{}}
\LANGes{
\Question{
Identifica cuál de las gráficas representa a la función \( f(x)=x+1-|x+1|;\ x\in [-4;4] \).}
{}{}{}{}{}{}}

\IMGanswer{1103072506a_0.pdf}
\IMGanswer{1103072506b_0.pdf}
\IMGanswer{1103072506c_0.pdf}
\IMGanswer{1103072506d_0.pdf}\End

\Data\ID{1103162908}
\Updated{2021-04-11 09:04:07}
\Level{B}
\SubArea{Functions with absolute values}
\IMG{1103162908.pdf}
\LANGen{
\Question{
With aid of the given graph of a function \( f \), find all \( x \) so that \( |f(x)-2|=1 \).}
{\( x\in\{-4;-2\} \)}{\( x\in\{-4;2\} \)}{\( x\in\{-2;2\} \)}{\( x\in\{-3\} \)}{}{}}
\LANGpl{
\Question{
Mając podany wykres funkcji \( f \), wyznacz zbiór rozwiązań równania \( |f(x)-2|=1 \).}
{\(x\in \{-4;-2\} \)}{\(x\in \{-4;2\} \)}{\(x\in \{-2;2\} \)}{\(x\in \{-3\} \)}{}{}}
\LANGcs{
\Question{
Pomocí grafu funkce \( f \) určete všechna \( x\), pro která platí: \( |f(x)-2|=1 \).}
{\(x\in \{-4;-2\} \)}{\(x\in \{-4;2\} \)}{\(x\in \{-2;2\} \)}{\(x\in \{-3\} \)}{}{}}
\LANGsk{
\Question{
Pomocou grafu funkcie \( f \) určte všetky \( x\), pre ktoré platí:  \( |f(x)-2|=1 \).}
{\( x\in\{-4;-2\} \)}{\(x\in \{-4;2\} \)}{\( x\in\{-2;2\} \)}{\( x\in\{-3\} \)}{}{}}
\LANGes{
\Question{
Con ayuda de la gráfica de la función \( f \), encuentra todos los \( x \) tales que \( |f(x)-2|=1 \).}
{\( x\in\{-4;-2\} \)}{\( x\in\{-4;2\} \)}{\( x\in\{-2;2\} \)}{\( x\in\{-3\} \)}{}{}}
\End

\Data\ID{1103162909}
\Updated{2020-11-24 11:11:03}
\Level{B}
\SubArea{Functions with absolute values}
\IMG{1103162909.pdf}
\LANGen{
\Question{
With aid of the graph of a function \( f \) find all \( x\), so that \( |f(x)|=3 \).}
{\( x\in\{-5;1\} \)}{\(x\in \{1\} \)}{\(x\in \{-2\} \)}{\( x\in\{-5;5\} \)}{}{}}
\LANGpl{
\Question{
Dany jest wykres funkcji \( f \), wyznacz zbiór rozwiązań równania \( |f(x)|=3 \).}
{\(x\in \{-5;1\} \)}{\(x\in \{1\} \)}{\(x\in \{-2\} \)}{\(x\in \{-5;5\} \)}{}{}}
\LANGcs{
\Question{
Pomocí grafu funkce \( f \) určete všechna \( x\), pro která platí:  \( |f(x)|=3 \).}
{\( x\in\{-5;1\} \)}{\(x\in \{1\} \)}{\( x\in\{-2\} \)}{\( x\in\{-5;5\} \)}{}{}}
\LANGsk{
\Question{
Pomocou grafu funkcie \( f \)  určte všetky \( x\), pre ktoré platí: \( |f(x)|=3 \).}
{\(x\in \{-5;1\} \)}{\(x\in \{1\} \)}{\(x\in \{-2\} \)}{\( x\in\{-5;5\} \)}{}{}}
\LANGes{
\Question{
Con ayuda de la gráfica de la función \( f \) encuentra todos \( x\), tales que \( |f(x)|=3 \).}
{\( x\in\{-5;1\} \)}{\(x\in \{1\} \)}{\(x\in \{-2\} \)}{\( x\in\{-5;5\} \)}{}{}}
\End

\Data\ID{1003049202}
\Updated{2022-04-23 05:04:37}
\Level{C}
\SubArea{Functions with absolute values}
\LANGen{
\Question{
Let \( f(x)=|x|-|3-2|x| | \). Identify which of the following function values is the smallest.}
{\( f(-6.5) \)}{\( f(-2) \)}{\( f(0) \)}{\( f(0.5) \)}{}{}}
\LANGpl{
\Question{
Niech \( f(x)=|x|-|3-2|x| | \). Określ, która z podanych wartości funkcji jest najniższa.}
{\( f(-6{,}5) \)}{\( f(-2) \)}{\( f(0) \)}{\( f(0{,}5) \)}{}{}}
\LANGcs{
\Question{
Funkce f je dána předpisem \( f(x)=|x|-|3-2|x| | \). Která z daných funkčních hodnot je nejmenší?}
{\( f(-6{,}5) \)}{\( f(-2) \)}{\( f(0) \)}{\( f(0{,}5) \)}{}{}}
\LANGsk{
\Question{
Daná je funkcia \( f(x)=|x|-|3-2|x| | \). Ktorá z nasledujúcich funkčných hodnôt je najmenšia?}
{\( f(-6{,}5) \)}{\( f(-2) \)}{\( f(0) \)}{\( f(0{,}5) \)}{}{}}
\LANGes{
\Question{
Sea \( f(x)=|x|-|3-2|x| | \). Identifica cuál de los valores funcionales es el más pequeño.}
{\( f(-6.5) \)}{\( f(-2) \)}{\( f(0) \)}{\( f(0.5) \)}{}{}}
\End

\Data\ID{1003049204}
\Updated{2020-11-24 08:11:22}
\Level{C}
\SubArea{Functions with absolute values}
\LANGen{
\Question{
Let \( f(x)=|x| \). Identify which of the statements is false.}
{\( \forall a\text{, }b\in\mathbb{R}\colon f(a+b)=f(a)+f(b) \)}{\( \forall a\text{, }b\in\mathbb{R}\colon f(a\cdot b)=f(a)\cdot f(b) \)}{\( \forall a\in\mathbb{R}\text{, }b\in\mathbb{R}\setminus\{0\}\colon f(\frac ab)=\frac{f(a)}{f(b)}  \)}{\( \forall a\in\mathbb{R}\colon f(a)=f(-a) \)}{}{}}
\LANGpl{
\Question{
Niech \( f(x)=|x| \). Określ, które z poniższych stwierdzeń jest  fałszywe.}
{\( \forall a\text{, }b\in\mathbb{R}\colon f(a+b)=f(a)+f(b) \)}{\( \forall a\text{, }b\in\mathbb{R}\colon f(a\cdot b)=f(a)\cdot f(b) \)}{\( \forall a\in\mathbb{R}\text{, }b\in\mathbb{R}\setminus\{0\}\colon f(\frac ab)=\frac{f(a)}{f(b)}  \)}{\( \forall a\in\mathbb{R}\colon f(a)=f(-a) \)}{}{}}
\LANGcs{
\Question{
Funkce f je dána předpisem \( f(x)=|x| \). Vyberte nepravdivý výrok:}
{\( \forall a\text{, }b\in\mathbb{R}\colon f(a+b)=f(a)+f(b) \)}{\( \forall a\text{, }b\in\mathbb{R}\colon f(a\cdot b)=f(a)\cdot f(b) \)}{\( \forall a\in\mathbb{R}\text{, }b\in\mathbb{R}\setminus\{0\}\colon f(\frac ab)=\frac{f(a)}{f(b)}  \)}{\( \forall a\in\mathbb{R}\colon f(a)=f(-a) \)}{}{}}
\LANGsk{
\Question{
Daná je funkcia \( f(x)=|x| \). Ktoré tvrdenie nie je pravdivé?}
{\( \forall a\text{, }b\in\mathbb{R}\colon f(a+b)=f(a)+f(b) \)}{\( \forall a\text{, }b\in\mathbb{R}\colon f(a\cdot b)=f(a)\cdot f(b) \)}{\( \forall a\in\mathbb{R}\text{, }b\in\mathbb{R}\setminus\{0\}\colon f(\frac ab)=\frac{f(a)}{f(b)}  \)}{\( \forall a\in\mathbb{R}\colon f(a)=f(-a) \)}{}{}}
\LANGes{
\Question{
Sea \( f(x)=|x| \). Identifica cuál de las declaraciones es falsa.}
{\( \forall a\text{, }b\in\mathbb{R}\colon f(a+b)=f(a)+f(b) \)}{\( \forall a\text{, }b\in\mathbb{R}\colon f(a\cdot b)=f(a)\cdot f(b) \)}{\( \forall a\in\mathbb{R}\text{, }b\in\mathbb{R}\setminus\{0\}\colon f(\frac ab)=\frac{f(a)}{f(b)}  \)}{\( \forall a\in\mathbb{R}\colon f(a)=f(-a) \)}{}{}}
\End

\Data\ID{1003102302}
\Updated{2021-04-11 09:04:26}
\Level{C}
\SubArea{Functions with absolute values}
\LANGen{
\Question{
Let \( f(x)=|1-|x| | \). Which of the following statements is true?}
{Function \( f \) has minimum at \( x=-1 \).}{Function \( f \) is bounded.}{Function \( f \) is increasing on the interval \( (0;\infty) \).}{The range of \( f \) is \( [1;\infty) \).}{}{}}
\LANGpl{
\Question{
Załóżmy, że \( f(x)=|1-|x| | \). Które z poniższych stwierdzeń jest prawdziwe?}
{Funkcja \( f \) osiąga minimum w \( x=-1 \).}{Funkcja \( f \) jest ograniczona.}{Funkcja \( f \) jest rosnąca w przedziale  \( (0;\infty) \).}{Zakresem funkcji  \( f \) jest \( \langle1;\infty) \).}{}{}}
\LANGcs{
\Question{
Funkce \( f \) je dána předpisem \( f(x)=|1-|x| | \). Vyberte pravdivý výrok.}
{Funkce \( f \) má minimum v bodě \( x=-1 \).}{Funkce \( f \) je omezená.}{Funkce \( f \) je rostoucí v intervalu \( (0;\infty) \).}{Obor hodnot funkce \( f \) je \( \langle1;\infty) \).}{}{}}
\LANGsk{
\Question{
Daná je funkcia predpisom \( f(x)=|1-|x| | \). Vyberte pravdivé tvrdenie o funkcii \( f \).}
{Funkcia \( f \) má minimum v bode \( x=-1 \).}{Funkcia \( f \) je ohraničená.}{Funkcia \( f \) je rastúca na intervale \( (0;\infty) \).}{Obor hodnôt funkcie \( f \) je interval \( \langle1;\infty) \).}{}{}}
\LANGes{
\Question{
Sea \( f(x)=|1-|x| | \). ¿Cuál de las declaraciones es correcta?}
{La función \( f \) tiene un mínimo en \( x=-1 \).}{La función \( f \) es acotada.}{La función \( f \) es creciente en el intervalo \( (0;\infty) \).}{El Rango de \( f \) es \( [1;\infty) \).}{}{}}
\End

\Data\ID{1003102303}
\Updated{2020-11-24 09:11:45}
\Level{C}
\SubArea{Functions with absolute values}
\LANGen{
\Question{
Which of the following functions is decreasing on the interval  \( (-\infty;0) \)?}
{\( m(x)=\left|x-|x-1|\right| \)}{\( h(x)=\left| x+|x-1|\right| \)}{\( g(x)=\left|x-|x+1|\right| \)}{\( f(x)=\left|x+|x+1|\right| \)}{}{}}
\LANGpl{
\Question{
Która z poniższych funkcji jest malejąca w przedziale \( (-\infty;0) \)?}
{\( m(x)=\left|x-|x-1|\right| \)}{\( h(x)=\left| x+|x-1|\right| \)}{\( g(x)=\left|x-|x+1|\right| \)}{\( f(x)=\left|x+|x+1|\right| \)}{}{}}
\LANGcs{
\Question{
Která z následujících funkcí je klesající v intervalu  \( (-\infty;0) \)?}
{\( m(x)=\left|x-|x-1|\right| \)}{\( h(x)=\left| x+|x-1|\right| \)}{\( g(x)=\left|x-|x+1|\right| \)}{\( f(x)=\left|x+|x+1|\right| \)}{}{}}
\LANGsk{
\Question{
Ktorá z nasledujúcich funkcií je klesajúca na intervale \( (-\infty;0) \)?}
{\( m(x)=\left|x-|x-1|\right| \)}{\( h(x)=\left| x+|x-1|\right| \)}{\( g(x)=\left|x-|x+1|\right| \)}{\( f(x)=\left|x+|x+1|\right| \)}{}{}}
\LANGes{
\Question{
¿Cuál de las siguientes funciones es decreciente en el intervalo  \( (-\infty;0) \)?}
{\( m(x)=\left|x-|x-1|\right| \)}{\( h(x)=\left| x+|x-1|\right| \)}{\( g(x)=\left|x-|x+1|\right| \)}{\( f(x)=\left|x+|x+1|\right| \)}{}{}}
\End

\Data\ID{1003187205}
\Updated{2021-04-11 09:04:59}
\Level{C}
\SubArea{Functions with absolute values}
\LANGen{
\Question{
Let \( f(x)=\left|3|2x-1|-9\right| \). The number of argument values \( x \) for which \( f(x)=2 \) is:}
{\( 4 \)}{\( 2 \)}{\( 3 \)}{\( 1 \)}{}{}}
\LANGpl{
\Question{
Niech  \( f(x)=\left|3|2x-1|-9\right| \).  Liczba wartości argumentu \( x \) dla którego \( f(x)=2 \) wynosi:}
{\( 4 \)}{\( 2 \)}{\( 3 \)}{\( 1 \)}{}{}}
\LANGcs{
\Question{
Nechť \( f(x)=\left|3|2x-1|-9\right| \). Počet hodnot \( x \), pro které platí \( f(x)=2 \), je:}
{\( 4 \)}{\( 2 \)}{\( 3 \)}{\( 1 \)}{}{}}
\LANGsk{
\Question{
Daná je funkcia \( f(x)=\left|3|2x-1|-9\right| \).  Počet hodnôt $x$, pre ktoré platí \( f(x)=2 \), je:}
{\( 4 \)}{\( 2 \)}{\( 3 \)}{\( 1 \)}{}{}}
\LANGes{
\Question{
Sea \( f(x)=\left|3|2x-1|-9\right| \). El valor de la variable \( x \) para el cuál \( f(x)=2 \) es:}
{\( 4 \)}{\( 2 \)}{\( 3 \)}{\( 1 \)}{}{}}
\End

\Data\ID{1003187206}
\Updated{2020-11-24 09:11:45}
\Level{C}
\SubArea{Functions with absolute values}
\LANGen{
\Question{
How many \( x \)-intercepts does the graph of \( f(x)=\left|-|2-x|- 2\right| \) have?}
{\( 0 \)}{\( 2 \)}{\( 1 \)}{\( 4 \)}{}{}}
\LANGpl{
\Question{
Ile punktów przecięcia z osią współrzędnych \( x \) posiada wykres funkcji \( f(x)=\left|-|2-x|- 2\right| \)?}
{\( 0 \)}{\( 2 \)}{\( 1 \)}{\( 4 \)}{}{}}
\LANGcs{
\Question{
Kolik průsečíku s osou \( x \) má graf funkce \( f(x)=\left|-|2-x|- 2\right| \)?}
{\( 0 \)}{\( 2 \)}{\( 1 \)}{\( 4 \)}{}{}}
\LANGsk{
\Question{
Koľko priesečníkov s osou \( x \) má graf funkcie \( f(x)=\left|-|2-x|- 2\right| \)?}
{\( 0 \)}{\( 2 \)}{\( 1 \)}{\( 4 \)}{}{}}
\LANGes{
\Question{
¿Cuántas veces interseca la gráfica de la función \( f(x)=\left|-|2-x|- 2\right| \) al eje \( x \)?}
{\( 0 \)}{\( 2 \)}{\( 1 \)}{\( 4 \)}{}{}}
\End

\Data\ID{1103102301}
\Updated{2021-04-11 09:04:47}
\Level{C}
\SubArea{Functions with absolute values}
\LANGen{
\Question{
Let \( f(x)=|x-|x| | \), \( x\in [-5;5] \). Choose the graph of \( f \).}
{}{}{}{}{}{}}
\LANGpl{
\Question{
Załóżmy, że \( f(x)=|x-|x| | \), \( x\in\langle-5;5\rangle \). Wybierz wykres funkcji  \( f \).}
{}{}{}{}{}{}}
\LANGcs{
\Question{
Na kterém obrázku je graf funkce  \( f(x)=|x-|x| | \), \( x\in\langle-5;5\rangle \)?}
{}{}{}{}{}{}}
\LANGsk{
\Question{
Ktorý z následujúcich grafov je graf funkcie \( f(x)=|x-|x| | \), \( x\in\langle-5;5\rangle \)?}
{}{}{}{}{}{}}
\LANGes{
\Question{
Sea \( f(x)=|x-|x| | \), \( x\in [-5;5] \). Elige la gráfica de \( f \).}
{}{}{}{}{}{}}

\IMGanswer{1103102301a.pdf}
\IMGanswer{1103102301b_0.pdf}
\IMGanswer{1103102301c_0.pdf}
\IMGanswer{1103102301d_0.pdf}\End

\Data\ID{2100002001}
\Updated{2021-12-26 08:12:12}
\Level{C}
\SubArea{Functions with absolute values}
\LANGen{
\Question{
Identify a possible graph of the function \(f(x) = \frac{1}{|x|}-1\).}
{}{}{}{}{}{}}
\LANGpl{
\Question{
Wskaż wykres funkcji  \(f(x) = \frac{1}{|x|}-1\).}
{}{}{}{}{}{}}
\LANGcs{
\Question{
Určete, který z grafů může být grafem funkce dané předpisem \(f(x) = \frac{1}{|x|}-1\).}
{}{}{}{}{}{}}
\LANGsk{
\Question{
Určte, ktorý z grafov môže byť grafom funkcie danej predpisom \(f(x) = \frac{1}{|x|}-1\).}
{}{}{}{}{}{}}
\LANGes{
\Question{
Identifica la gráfica de la función \(f(x) = \frac{1}{|x|}-1\).}
{}{}{}{}{}{}}

\IMGanswer{linlomabs-figure0.pdf}
\IMGanswer{linlomabs-figure1.pdf}
\IMGanswer{linlomabs-figure2.pdf}
\IMGanswer{linlomabs-figure3.pdf}\End

\Data\ID{2100002002}
\Updated{2021-12-26 08:12:27}
\Level{C}
\SubArea{Functions with absolute values}
\LANGen{
\Question{
Identify a possible graph of the function \(f(x) = \frac{1}{|x-1|}+1\).}
{}{}{}{}{}{}}
\LANGpl{
\Question{
Wskaż wykres funkcji \(f(x) = \frac{1}{|x-1|}+1\).}
{}{}{}{}{}{}}
\LANGcs{
\Question{
Určete, který z grafů může být grafem funkce dané předpisem \(f(x) = \frac{1}{|x-1|}+1\).}
{}{}{}{}{}{}}
\LANGsk{
\Question{
Určte, ktorý z grafov môže byť grafom funkcie danej predpisom \(f(x) = \frac{1}{|x-1|}+1\).}
{}{}{}{}{}{}}
\LANGes{
\Question{
Identifica la gráfica de la función \(f(x) = \frac{1}{|x-1|}+1\).}
{}{}{}{}{}{}}

\IMGanswer{linlomabs-figure4.pdf}
\IMGanswer{linlomabs-figure5.pdf}
\IMGanswer{linlomabs-figure6.pdf}
\IMGanswer{linlomabs-figure7.pdf}\End

\Data\ID{2100002003}
\Updated{2021-12-26 08:12:46}
\Level{C}
\SubArea{Functions with absolute values}
\LANGen{
\Question{
Identify a possible graph of the function \(f(x) = \frac{1}{|x+1|}-1\).}
{}{}{}{}{}{}}
\LANGpl{
\Question{
Wskaż wykres funkcji \(f(x) = \frac{1}{|x+1|}-1\).}
{}{}{}{}{}{}}
\LANGcs{
\Question{
Určete, který z grafů může být grafem funkce dané předpisem \(f(x) = \frac{1}{|x+1|}-1\).}
{}{}{}{}{}{}}
\LANGsk{
\Question{
Určte, ktorý z grafov môže byť grafom funkcie danej predpisom \(f(x) = \frac{1}{|x+1|}-1\).}
{}{}{}{}{}{}}
\LANGes{
\Question{
Identifica la gráfica de la función \(f(x) = \frac{1}{|x+1|}-1\).}
{}{}{}{}{}{}}

\IMGanswer{linlomabs-figure8.pdf}
\IMGanswer{linlomabs-figure9.pdf}
\IMGanswer{linlomabs-figure10.pdf}
\IMGanswer{linlomabs-figure11.pdf}\End

\Data\ID{2100002004}
\Updated{2021-12-26 08:12:07}
\Level{C}
\SubArea{Functions with absolute values}
\LANGen{
\Question{
Identify a possible graph of the function \( f(x) =\left| \frac{1}{x+1} -1 \right| \).}
{}{}{}{}{}{}}
\LANGpl{
\Question{
Wskaż wykres funkcji \( f(x) =\left| \frac{1}{x+1} -1 \right| \).}
{}{}{}{}{}{}}
\LANGcs{
\Question{
Určete, který z grafů může být grafem funkce dané předpisem \( f(x) =\left| \frac{1}{x+1} -1 \right| \).}
{}{}{}{}{}{}}
\LANGsk{
\Question{
Určte, ktorý z grafov môže byť grafom funkcie danej predpisom \( f(x) =\left| \frac{1}{x+1} -1 \right| \).}
{}{}{}{}{}{}}
\LANGes{
\Question{
Identifica la gráfica de la función  \( f(x) =\left| \frac{1}{x+1} -1 \right| \).}
{}{}{}{}{}{}}

\IMGanswer{linlomabs-figure12.pdf}
\IMGanswer{linlomabs-figure13.pdf}
\IMGanswer{linlomabs-figure14.pdf}
\IMGanswer{linlomabs-figure15.pdf}\End

\Data\ID{2100002005}
\Updated{2021-12-26 08:12:26}
\Level{C}
\SubArea{Functions with absolute values}
\LANGen{
\Question{
Identify a possible graph of the function \(f(x) = \left| \frac{1}{x-1} +1\right|\).}
{}{}{}{}{}{}}
\LANGpl{
\Question{
Wskaż wykres funkcji \(f(x) = \left| \frac{1}{x-1} +1\right|\).}
{}{}{}{}{}{}}
\LANGcs{
\Question{
Určete, který z grafů může být grafem funkce dané předpisem \(f(x) = \left| \frac{1}{x-1} +1\right|\).}
{}{}{}{}{}{}}
\LANGsk{
\Question{
Určte, ktorý z grafov môže byť grafom funkcie danej predpisom \(f(x) = \left| \frac{1}{x-1} +1\right|\).}
{}{}{}{}{}{}}
\LANGes{
\Question{
Identifica la gráfica de la función \(f(x) = \left| \frac{1}{x-1} +1\right|\).}
{}{}{}{}{}{}}

\IMGanswer{linlomabs-figure16.pdf}
\IMGanswer{linlomabs-figure17.pdf}
\IMGanswer{linlomabs-figure18.pdf}
\IMGanswer{linlomabs-figure19.pdf}\End

\Data\ID{9100004009}
\Updated{2020-11-24 02:11:18}
\Level{C}
\SubArea{Functions with absolute values}
\LANGen{
\Question{
Identify a possible graph of the function
\(g(x) ={\Bigl | |x - 2|- 1\Bigr |}\).}
{}{}{}{}{}{}}
\LANGpl{
\Question{
Określ, który z rysunków przedstawia wykres funkcji
\(g\colon y ={\Bigl | |x - 2|- 1\Bigr |}\).}
{}{}{}{}{}{}}
\LANGcs{
\Question{
Určete, který z obrázků je grafem funkce
\(g\colon y ={\Bigl | |x - 2|- 1\Bigr |}\).}
{}{}{}{}{}{}}
\LANGsk{
\Question{
Ktorý z grafov je grafom funkcie
\(g\colon y ={\Bigl | |x - 2|- 1\Bigr |}\)?}
{}{}{}{}{}{}}
\LANGes{
\Question{
Identifica la posible gráfica de la función
\(g(x) ={\Bigl | |x - 2|- 1\Bigr |}\).}
{}{}{}{}{}{}}

\IMGanswer{000040_09-figure3.pdf}
\IMGanswer{000040_09-figure0.pdf}
\IMGanswer{000040_09-figure1.pdf}
\IMGanswer{000040_09-figure2.pdf}\End

\Data\ID{9100004010}
\Updated{2020-11-24 02:11:30}
\Level{C}
\SubArea{Functions with absolute values}
\LANGen{
\Question{
Identify a possible graph of the function
\(g\colon y ={\Bigl | |x - 1|- 1\Bigr |} - 1\).}
{}{}{}{}{}{}}
\LANGpl{
\Question{
Określ, który z poniższych rysunków przedstawia wykres funkcji
\(g\colon y ={\Bigl | |x - 1|- 1\Bigr |} - 1\).}
{}{}{}{}{}{}}
\LANGcs{
\Question{
Určete, který z obrázků je grafem funkce
\(g\colon y ={\Bigl | |x - 1|- 1\Bigr |} - 1\).}
{}{}{}{}{}{}}
\LANGsk{
\Question{
Ktorý z nasledujúcich grafov je grafom funkcie
\(g\colon y ={\Bigl | |x - 1|- 1\Bigr |} - 1\)?}
{}{}{}{}{}{}}
\LANGes{
\Question{
Identifica la posible gráfica de la función
\(g\colon y ={\Bigl | |x - 1|- 1\Bigr |} - 1\).}
{}{}{}{}{}{}}

\IMGanswer{000040_10-figure0.pdf}
\IMGanswer{000040_10-figure1.pdf}
\IMGanswer{000040_10-figure2.pdf}
\IMGanswer{000040_10-figure3.pdf}\End
